% Plantilla para un Trabajo Fin de Grado de la Universidad de Granada,
% adaptada para el Doble Grado en Ingeniería Informática y Matemáticas.
%
%  Autor: Mario Román.
%  Licencia: GNU GPLv2.
%
% Esta plantilla es una adaptación al castellano de la plantilla
% classicthesis de André Miede, que puede obtenerse en:
%  https://ctan.org/tex-archive/macros/latex/contrib/classicthesis?lang=en
% La plantilla original se licencia en GNU GPLv2.
%
% Esta plantilla usa símbolos de la Universidad de Granada sujetos a la normativa
% de identidad visual corporativa, que puede encontrarse en:
% http://secretariageneral.ugr.es/pages/ivc/normativa
%
% La compilación se realiza con las siguientes instrucciones:
%   pdflatex --shell-escape main.tex
%   bibtex main
%   pdflatex --shell-escape main.tex
%   pdflatex --shell-escape main.tex

% Opciones del tipo de documento
\documentclass[twoside,openright,titlepage,numbers=noenddot,headinclude,footinclude=true,
cleardoublepage=empty,abstractoff,BCOR=5mm,paper=a4,fontsize=12pt,main=spanish,pdfstartpage=1]{scrreprt}

% Paquetes de latex que se cargan al inicio. Cubren la entrada de
% texto, gráficos, código fuente y símbolos.
\PassOptionsToPackage{table,dvipsnames}{xcolor}

\usepackage[utf8]{inputenc}
\usepackage[T1]{fontenc}
\usepackage{fixltx2e}
\usepackage{graphicx} % Inclusión de imágenes.
\usepackage{grffile}  % Distintos formatos para imágenes.
\usepackage{longtable} % Tablas multipágina.
\usepackage{wrapfig} % Coloca texto alrededor de una figura.
\usepackage{rotating}
\usepackage[normalem]{ulem}
\usepackage{amsmath}
\usepackage{dsfont}
\usepackage{textcomp}
\usepackage{amssymb}
\usepackage{capt-of}
\usepackage[colorlinks=true]{hyperref}
\usepackage{tikz} % Diagramas conmutativos.
\usepackage{minted} % Código fuente.
\usepackage[numbers]{natbib}
\usepackage{booktabs} % Para tablas de figuras (ejemplo en salesman)
\usepackage{listings} % Código
\usepackage{comment}
\usepackage{float}
\usepackage[table]{xcolor}
\PassOptionsToPackage{table,dvipsnames}{xcolor}

\lstset{
	basicstyle=\ttfamily,%
	breaklines=true,%
	captionpos=b,                    % sets the caption-position to bottom
  tabsize=2,	                   % sets default tabsize to 2 spaces
  frame=none,
  numbers=left,
  xleftmargin=18pt,
  stepnumber=1,
  aboveskip=12pt,
  showstringspaces=false,
  keywordstyle=\bfseries,
  commentstyle=\itshape,
  numberstyle=\scriptsize\bfseries,
  %morekeywords={sage},
}


% Operadores

\DeclareMathOperator{\gr}{gr}
\DeclareMathOperator{\mcd}{mcd}

\lstset{ frame=Ltb,
framerule=0pt,
aboveskip=0.5cm,
framextopmargin=3pt,
framexbottommargin=3pt,
framexleftmargin=0.4cm,
framesep=0pt,
rulesep=.4pt,
backgroundcolor=\color{white},
rulesepcolor=\color{black},
%
stringstyle=\ttfamily,
showstringspaces = false,
basicstyle=\small\ttfamily,
commentstyle=\color{gray},
keywordstyle=\bfseries,
%
numbers=left,
numbersep=15pt,
numberstyle=\tiny,
numberfirstline = false,
breaklines=true,
}


\definecolor{codegreen}{rgb}{0,0.6,0}
\definecolor{codegray}{rgb}{0.5,0.5,0.5}
\definecolor{codepurple}{rgb}{0.58,0,0.82}
\definecolor{backcolour}{rgb}{0.95,0.95,0.92}


\lstdefinestyle{mystyle}{
  language=C++,
    commentstyle=\color{codegreen},
    keywordstyle=\color{magenta},
    numberstyle=\tiny\color{codegray},
    stringstyle=\color{codepurple},
    basicstyle=\ttfamily\footnotesize,
}

\lstset{style=mystyle}
% Plantilla classicthesis
\usepackage[beramono,eulerchapternumbers,linedheaders,parts,a5paper,dottedtoc,
manychapters,pdfspacing]{classicthesis}

% Geometría y espaciado de párrafos.
\setcounter{secnumdepth}{0}
\usepackage{enumitem}
\setitemize{noitemsep,topsep=0pt,parsep=0pt,partopsep=0pt}
\setlist[enumerate]{topsep=0pt,itemsep=-1ex,partopsep=1ex,parsep=1ex}
\usepackage[top=1in, bottom=1.5in, left=1in, right=1in]{geometry}
\setlength\itemsep{0em}
\setlength{\parindent}{0pt}
\usepackage{parskip}
\usepackage{setspace}

% Profundidad de la tabla de contenidos.
\setcounter{secnumdepth}{3}

% Usa el paquete minted para mostrar trozos de código.
% Pueden seleccionarse el lenguaje apropiado y el estilo del código.
\usepackage{minted}
\usemintedstyle{colorful}
\setminted{fontsize=\small}
\setminted[haskell]{linenos=false,fontsize=\small}
\renewcommand{\theFancyVerbLine}{\sffamily\textcolor[rgb]{0.5,0.5,1.0}{\oldstylenums{\arabic{FancyVerbLine}}}}

% Path para las imágenes
\graphicspath{{figures/}}

% Archivos de configuración.
%------------------------
% Bibliotecas para matemáticas de latex
%------------------------
\usepackage{amsthm}
\usepackage{amsmath}
\usepackage{tikz}
\usepackage{tikz-cd}
\usetikzlibrary{shapes,fit}
\usepackage{bussproofs}
\EnableBpAbbreviations{}
\usepackage{mathtools}
\usepackage{scalerel}
\usepackage{stmaryrd}


%------------------------
% Estilos para los teoremas
%------------------------
\theoremstyle{plain}
\newtheorem{theorem}{Teorema}
\newtheorem{proposition}{Proposición}
\newtheorem{lemma}{Lema}
\newtheorem{corollary}{Corolario}

\theoremstyle{definition}
\newtheorem{definition}{Definición}

% Change the proof style so it's in English and add \qed at the end.
\renewenvironment{proof}{{\bfseries Demostración.}}{\qed}

\theoremstyle{remark}
\newtheorem{remark}{Observación}
\newtheorem{exampleth}{Ejemplo}

% New style for postulates so they are tabulated
\makeatletter
\newtheoremstyle{indented}
	{3pt}% space before
	{3pt}% space after
	{\addtolength{\@totalleftmargin}{3.5em}
		\addtolength{\linewidth}{-3.5em}
		\parshape 1 3.5em \linewidth}% body font
	{}% indent
	{\bfseries}% header font
	{.}% punctuation
	{.5em}% after theorem header
	{}% header specification (empty for default)
\makeatother

% Apply the new style
\theoremstyle{indented}
\newtheorem{postulate}{Postulate}
\newtheorem*{postulate 3'}{Postulate 3'}
\newtheorem*{postulate 2'}{Projective Measurement}

%------------------------
% Macros
% ------------------------

\newcommand*{\B}{\mathbb{B}}
\newcommand*{\C}{\mathds{C}}
\newcommand*{\R}{\mathbb{R}}
\newcommand*{\ra}{\rangle}
\newcommand*{\la}{\langle}

% Para poner sonrisa sobre puntos suspensivos. Uso: \overplace{n}{\dotsc}
\newcommand{\overplace}[2]{%
	\overset{\substack{#1\\\smile}}{#2}%
}

% Para añadir un salto de línea al nuevo párrafo. Use: \newparagraph
\newcommand{\newparagraph}[1]{\paragraph{#1}\mbox{}\\}  % En macros.tex se almacenan las opciones y comandos para escribir matemáticas.
% ****************************************************************************************************
% classicthesis-config.tex 
% formerly known as loadpackages.sty, classicthesis-ldpkg.sty, and classicthesis-preamble.sty 
% Use it at the beginning of your ClassicThesis.tex, or as a LaTeX Preamble 
% in your ClassicThesis.{tex,lyx} with \input{classicthesis-config}
% ****************************************************************************************************  
% If you like the classicthesis, then I would appreciate a postcard. 
% My address can be found in the file ClassicThesis.pdf. A collection 
% of the postcards I received so far is available online at 
% http://postcards.miede.de
% ****************************************************************************************************


% ****************************************************************************************************
% 0. Set the encoding of your files. UTF-8 is the only sensible encoding nowadays. If you can't read
% äöüßáéçèê∂åëæƒÏ€ then change the encoding setting in your editor, not the line below. If your editor
% does not support utf8 use another editor!
% ****************************************************************************************************
\PassOptionsToPackage{utf8x}{inputenc}
	\usepackage{inputenc}

% ****************************************************************************************************
% 1. Configure classicthesis for your needs here, e.g., remove "drafting" below 
% in order to deactivate the time-stamp on the pages
% ****************************************************************************************************
\PassOptionsToPackage{eulerchapternumbers,listings,drafting,%
		pdfspacing,%floatperchapter,%linedheaders,%
                subfig,beramono,eulermath,parts,dottedtoc}{classicthesis}                                        
% ********************************************************************
% Available options for classicthesis.sty 
% (see ClassicThesis.pdf for more information):
% drafting
% parts nochapters linedheaders
% eulerchapternumbers beramono eulermath pdfspacing minionprospacing
% tocaligned dottedtoc manychapters
% listings floatperchapter subfig
% ********************************************************************

% ****************************************************************************************************
% 2. Personal data and user ad-hoc commands
% ****************************************************************************************************
\newcommand{\myTitle}{Prácticas DDSI\xspace}
\newcommand{\mySubtitle}{An Homage to The Elements of Typographic Style\xspace}
\newcommand{\myDegree}{Doktor-Ingenieur (Dr.-Ing.)\xspace}
\newcommand{\myName}{André Miede\xspace}
\newcommand{\myProf}{Put name here\xspace}
\newcommand{\myOtherProf}{Put name here\xspace}
\newcommand{\mySupervisor}{Put name here\xspace}
\newcommand{\myFaculty}{Put data here\xspace}
\newcommand{\myDepartment}{Put data here\xspace}
\newcommand{\myUni}{Put data here\xspace}
\newcommand{\myLocation}{Saarbrücken\xspace}
\newcommand{\myTime}{September 2015\xspace}
%\newcommand{\myVersion}{version 4.2\xspace}

% ********************************************************************
% Setup, finetuning, and useful commands
% ********************************************************************
\newcounter{dummy} % necessary for correct hyperlinks (to index, bib, etc.)
\newlength{\abcd} % for ab..z string length calculation
\providecommand{\mLyX}{L\kern-.1667em\lower.25em\hbox{Y}\kern-.125emX\@}
\newcommand{\ie}{i.\,e.}
\newcommand{\Ie}{I.\,e.}
\newcommand{\eg}{e.\,g.}
\newcommand{\Eg}{E.\,g.} 
% ****************************************************************************************************


% ****************************************************************************************************
% 3. Loading some handy packages
% ****************************************************************************************************
% ******************************************************************** 
% Packages with options that might require adjustments
% ******************************************************************** 
%\PassOptionsToPackage{ngerman,american}{babel}   % change this to your language(s)
% Spanish languages need extra options in order to work with this template
% \PassOptionsToPackage{es-lcroman,spanish}{babel}
\usepackage[main=spanish]{babel}


%\usepackage{csquotes}
% \PassOptionsToPackage{%
%     %backend=biber, %instead of bibtex
% 	backend=bibtex8,bibencoding=ascii,%
% 	language=auto,%
% 	style=alpha,%
%     %style=authoryear-comp, % Author 1999, 2010
%     %bibstyle=authoryear,dashed=false, % dashed: substitute rep. author with ---
%     sorting=nyt, % name, year, title
%     maxbibnames=10, % default: 3, et al.
%     %backref=true,%
%     natbib=true % natbib compatibility mode (\citep and \citet still work)
% }{biblatex}
%     \usepackage{biblatex}

% \PassOptionsToPackage{fleqn}{amsmath}       % math environments and more by the AMS 
%     \usepackage{amsmath}

% ******************************************************************** 
% General useful packages
% ******************************************************************** 
\PassOptionsToPackage{T1}{fontenc} % T2A for cyrillics
    \usepackage{fontenc}     
\usepackage{textcomp} % fix warning with missing font shapes
\usepackage{scrhack} % fix warnings when using KOMA with listings package          
\usepackage{xspace} % to get the spacing after macros right  
\usepackage{mparhack} % get marginpar right
\usepackage{fixltx2e} % fixes some LaTeX stuff --> since 2015 in the LaTeX kernel (see below)
%\usepackage[latest]{latexrelease} % will be used once available in more distributions (ISSUE #107)
\PassOptionsToPackage{printonlyused,smaller}{acronym} 
    \usepackage{acronym} % nice macros for handling all acronyms in the thesis
    %\renewcommand{\bflabel}[1]{{#1}\hfill} % fix the list of acronyms --> no longer working
    %\renewcommand*{\acsfont}[1]{\textsc{#1}} 
    \renewcommand*{\aclabelfont}[1]{\acsfont{#1}}
% ****************************************************************************************************


% ****************************************************************************************************
% 4. Setup floats: tables, (sub)figures, and captions
% ****************************************************************************************************
\usepackage{tabularx} % better tables
    \setlength{\extrarowheight}{3pt} % increase table row height
\newcommand{\tableheadline}[1]{\multicolumn{1}{c}{\spacedlowsmallcaps{#1}}}
\newcommand{\myfloatalign}{\centering} % to be used with each float for alignment
\usepackage{caption}
% Thanks to cgnieder and Claus Lahiri
% http://tex.stackexchange.com/questions/69349/spacedlowsmallcaps-in-caption-label
% [REMOVED DUE TO OTHER PROBLEMS, SEE ISSUE #82]    
%\DeclareCaptionLabelFormat{smallcaps}{\bothIfFirst{#1}{~}\MakeTextLowercase{\textsc{#2}}}
%\captionsetup{font=small,labelformat=smallcaps} % format=hang,
\captionsetup{font=small} % format=hang,
\usepackage{subfig}  
% ****************************************************************************************************


% ****************************************************************************************************
% 5. Setup code listings
% ****************************************************************************************************
% \usepackage{listings} 
% %\lstset{emph={trueIndex,root},emphstyle=\color{BlueViolet}}%\underbar} % for special keywords
% \lstset{language={Haskell},morekeywords={PassOptionsToPackage,selectlanguage},keywordstyle=\color{RoyalBlue},basicstyle=\small\ttfamily,commentstyle=\color{Green}\ttfamily,stringstyle=\rmfamily,numbers=none,numberstyle=\scriptsize,stepnumber=5,numbersep=8pt,showstringspaces=false,breaklines=true,belowcaptionskip=.75\baselineskip} 
% ****************************************************************************************************             


% ****************************************************************************************************
% 6. PDFLaTeX, hyperreferences and citation backreferences
% ****************************************************************************************************
% ********************************************************************
% Using PDFLaTeX
% ********************************************************************
\PassOptionsToPackage{pdftex,hyperfootnotes=false,pdfpagelabels}{hyperref}
    \usepackage{hyperref}  % backref linktocpage pagebackref
\pdfcompresslevel=9
\pdfadjustspacing=1 
\PassOptionsToPackage{pdftex}{graphicx}
    \usepackage{graphicx} 
 

% ********************************************************************
% Hyperreferences
% ********************************************************************
\hypersetup{%
    %draft, % = no hyperlinking at all (useful in b/w printouts)
    colorlinks=true, linktocpage=true, pdfstartpage=3, pdfstartview=FitV,%
    % uncomment the following line if you want to have black links (e.g., for printing)
    %colorlinks=false, linktocpage=false, pdfstartpage=3, pdfstartview=FitV, pdfborder={0 0 0},%
    breaklinks=true, pdfpagemode=UseNone, pageanchor=true, pdfpagemode=UseOutlines,%
    plainpages=false, bookmarksnumbered, bookmarksopen=true, bookmarksopenlevel=1,%
    hypertexnames=true, pdfhighlight=/O,%nesting=true,%frenchlinks,%
    urlcolor=webbrown, linkcolor=RoyalBlue, citecolor=webgreen, %pagecolor=RoyalBlue,%
    %urlcolor=Black, linkcolor=Black, citecolor=Black, %pagecolor=Black,%
    pdftitle={\myTitle},%
    pdfauthor={\textcopyright\ \myName, \myUni, \myFaculty},%
    pdfsubject={},%
    pdfkeywords={},%
    pdfcreator={pdfLaTeX},%
    pdfproducer={LaTeX with hyperref and classicthesis}%
}   

% ********************************************************************
% Setup autoreferences
% ********************************************************************
% There are some issues regarding autorefnames
% http://www.ureader.de/msg/136221647.aspx
% http://www.tex.ac.uk/cgi-bin/texfaq2html?label=latexwords
% you have to redefine the makros for the 
% language you use, e.g., american, ngerman
% (as chosen when loading babel/AtBeginDocument)
% ********************************************************************
\makeatletter
\@ifpackageloaded{babel}%
    {%
        
       \addto\extrasamerican{%
			\renewcommand*{\figureautorefname}{Figure}%
			\renewcommand*{\tableautorefname}{Table}%
			\renewcommand*{\partautorefname}{Part}%
			\renewcommand*{\chapterautorefname}{Chapter}%
			\renewcommand*{\sectionautorefname}{Section}%
			\renewcommand*{\subsectionautorefname}{Section}%
			\renewcommand*{\subsubsectionautorefname}{Section}%     
                }%
       \addto\extrasngerman{% 
			\renewcommand*{\paragraphautorefname}{Absatz}%
			\renewcommand*{\subparagraphautorefname}{Unterabsatz}%
			\renewcommand*{\footnoteautorefname}{Fu\"snote}%
			\renewcommand*{\FancyVerbLineautorefname}{Zeile}%
			\renewcommand*{\theoremautorefname}{Theorem}%
			\renewcommand*{\appendixautorefname}{Anhang}%
			\renewcommand*{\equationautorefname}{Gleichung}%        
			\renewcommand*{\itemautorefname}{Punkt}%
                }%  
            % Fix to getting autorefs for subfigures right (thanks to Belinda Vogt for changing the definition)
            \providecommand{\subfigureautorefname}{\figureautorefname}%             
    }{\relax}
\makeatother


% ****************************************************************************************************
% 7. Last calls before the bar closes
% ****************************************************************************************************
% ********************************************************************
% Development Stuff
% ********************************************************************
\listfiles
%\PassOptionsToPackage{l2tabu,orthodox,abort}{nag}
%   \usepackage{nag}
%\PassOptionsToPackage{warning, all}{onlyamsmath}
%   \usepackage{onlyamsmath}

% ********************************************************************
% Last, but not least...
% ********************************************************************
\usepackage{classicthesis} 
% ****************************************************************************************************


% ****************************************************************************************************
% 8. Further adjustments (experimental)
% ****************************************************************************************************
% ********************************************************************
% Changing the text area
% ********************************************************************
%\linespread{1.05} % a bit more for Palatino
%\areaset[current]{312pt}{761pt} % 686 (factor 2.2) + 33 head + 42 head \the\footskip
%\setlength{\marginparwidth}{7em}%
%\setlength{\marginparsep}{2em}%

% ********************************************************************
% Using different fonts
% ********************************************************************
%\usepackage[oldstylenums]{kpfonts} % oldstyle notextcomp
%\usepackage[osf]{libertine}
%\usepackage[light,condensed,math]{iwona}
%\renewcommand{\sfdefault}{iwona}
%\usepackage{lmodern} % <-- no osf support :-(
%\usepackage{cfr-lm} % 
%\usepackage[urw-garamond]{mathdesign} <-- no osf support :-(
%\usepackage[default,osfigures]{opensans} % scale=0.95 
%\usepackage[sfdefault]{FiraSans}
% **************************************************************************************************** % En classicthesis-config.tex se almacenan las opciones propias de la plantilla.

% Color institucional UGR
% \definecolor{ugrColor}{HTML}{ed1c3e} % Versión clara.
\definecolor{ugrColor}{HTML}{FFA500}  % Usado en el título.
\definecolor{ugrColor2}{HTML}{FFA500} % Usado en las secciones.
% Prácticas DDSI
% Definición de colores
\definecolor{rojo}{RGB}{200,0,0}
\definecolor{azul}{RGB}{0,70,160}
\definecolor{verde}{RGB}{0,160,70}

% Datos de portada
\usepackage{titling} % Facilita los datos de la portada
\author{Ismael Sallami Moreno \\ Jesús Rodríguez González \\ Fernando José Gracia Choin \\ Sergio Calvo González \\ Javier Niño Sánchez} 
\title{Lista estructurada de requisitos}

% Portada
\usepackage{datetime}
\renewcommand\maketitle{
  \begin{titlepage}
    \begin{addmargin}[-0cm]{-0cm}
      \begin{center}
        \normalsize  
        \hfill
        \vfill
        \vspace{1.8cm}

        \spacedallcaps{Diseño y Desarrollo de Sistemas de Información: Práctica 1} \\ \medskip 
        \small{\spacedallcaps{Doble grado en Ingeniería Informática + ADE}} \\  \bigskip\bigskip
        \vspace{1.5cm}

        \begingroup
        \LARGE{\color{ugrColor}\spacedallcaps{\textbf{\thetitle}}} \\ \bigskip
        \endgroup

        \vfill
        \vspace{1.5cm}

        \textbf{Autores}\\
        \spacedlowsmallcaps{\theauthor}

        \hfill
        \vspace{1.7cm}

        \quad
        \includegraphics{figures/etsiit_logo.png}

        \spacedlowsmallcaps{E.T.S. de Ingenierías Informática y de Telecomunicación} \\ \medskip
        \vspace{1.0cm}
        %\textit{Curso académico 2021-2022}
        \textit{Granada, a \today}

        \vfill                      

      \end{center}  
    \end{addmargin}       
  \end{titlepage}}
\usepackage{wallpaper}
\usepackage[main=spanish]{babel}


\usepackage{slantsc}

% Algoritmos

\usepackage[linesnumbered, onelanguage, ruled]{algorithm2e}
\SetAlCapFnt{\sffamily}
\SetAlCapNameFnt{\sffamily}
\newcommand\mycommentfont[1]{\sffamily\textcolor{darkgray}{#1}}
\SetCommentSty{mycommentfont}
%\SetAlFnt{\sffamily}
%\SetKwSty{sffamily}
\SetAlCapSkip{.5em}
\SetAlgoCaptionLayout{raggedright}
\SetAlgorithmName{Algoritmo}{Algoritmo}{Lista de algoritmos}
%\SetAlgoSkip{medskip}
\SetAlgoInsideSkip{medskip}

% Entorno algoritmo con líneas pero caption debajo
\makeatletter
\newenvironment{Ualgorithm}[1][htpb]{\def\@algocf@post@ruled{\color{gray}\hrule  height\algoheightrule\kern6pt\relax}%
\def\@algocf@capt@ruled{under}%
\setlength\algotitleheightrule{0pt}%
\SetAlgoCaptionLayout{raggedright}%
\begin{algorithm}[#1]}
{\end{algorithm}}
\makeatother

% Comandos
\theoremstyle{plain}
\newtheorem{problemth}{Problema}

%Tikz styles para la práctica 2 de ddsi
\usepackage{tikz}
\usetikzlibrary{
    calc,           % Para cálculos de coordenadas
    positioning,    % Para colocar nodos de manera relativa (ej: below of)
    shapes,         % Para formas avanzadas (diamantes, elipses, etc.)
    arrows,         % Para estilos de flechas
    fit             % Para ajustar un nodo a un conjunto de nodos
}

\tikzstyle{agente} = [
    rectangle, draw, thick, 
    minimum width=2.5cm, minimum height=1cm, 
    fill=blue!15, text centered, 
    font=\bfseries\small
] % Agente Externo

\tikzstyle{proceso} = [
    rectangle, rounded corners=5pt, draw, thick, 
    minimum width=2.5cm, minimum height=1cm, 
    fill=green!20, text centered, 
    font=\bfseries\small
] % Proceso

\tikzset{
    almacen/.style={
        rectangle,
        minimum width=3cm,
        minimum height=1cm,
        fill=yellow!20,
        text centered,
        inner sep=5pt,
        font=\bfseries\small,
        draw=none, % no dibuja borde automático
        append after command={
            \pgfextra{
                % Dibuja todos los lados excepto el derecho
                \draw[thick]
                    (\tikzlastnode.north west) --
                    (\tikzlastnode.north east) --
                    (\tikzlastnode.south east) --
                    (\tikzlastnode.south west) -- cycle;
                % Borra (no dibujamos) el lado derecho
                \draw[white, line width=3pt]
                    (\tikzlastnode.north east) --
                    (\tikzlastnode.south east);
            }
        }
    }
}


\tikzstyle{flujo} = [
  ->, very thick, >=stealth, 
  font=\small
] % Flujo de datos

\tikzset{
  flujo/.style={->, thick, shorten >=6pt, shorten <=6pt, to path={-| (\tikztotarget)}}
}

\tikzset{
  flujo/.style={
    ->, thick,
    shorten >=6pt, shorten <=6pt,
    rounded corners=8pt, % esquinas redondeadas para rodear nodos
  },
  autoavoid/.style={
    to path={
      % Intenta evitar otras flechas usando out/in y bend
      [rounded corners=12pt]
      |- (\tikztotarget)
      \tikztonodes
    }
  }
}

\tikzstyle{iflujo} = [
  <-, very thick, >=stealth, 
  font=\small
] % Flujo de datos inverso

\tikzset{
  iflujo/.style={<-, thick, shorten >=6pt, shorten <=6pt, to path={-| (\tikztotarget)}}
}

\tikzset{
  iflujo/.style={
    <-, thick,
    shorten >=6pt, shorten <=6pt,
    rounded corners=8pt, % esquinas redondeadas para rodear nodos
  },
  autoavoid/.style={
    to path={
      % Intenta evitar otras flechas usando out/in y bend
      [rounded corners=12pt]
      |- (\tikztotarget)
      \tikztonodes
    }
  }
}

\makeatletter
\@addtoreset{chapter}{part}
\makeatother




\begin{document}


\ThisULCornerWallPaper{1}{figures/ugrA4_orange.pdf}
\maketitle


\onehalfspacing

% Indice
\tableofcontents
\newpage



% \chapter{Ejemplo}
% \section{Ejemplo de requisito funcional}
\textbf{\textcolor{rojo}{\underline{RF1:}}} Dar de alta contacto. \\[6pt]

\textbf{\textcolor{rojo}{Entrada:}} Agente externo: usuario. Acción: solicitar inserción. 
Requisito de datos de entrada \textcolor{azul}{RDE1}. \\[6pt]

\textbf{\textcolor{rojo}{BD:}} Requisito de datos de escritura \textcolor{azul}{RDW1}. \\[6pt]

\textbf{\textcolor{rojo}{Salida:}} Agente externo: usuario. Acción: confirmación resultado. 
Requisito de datos de salida: ninguno. \\[12pt]

\textcolor{azul}{\textbf{RDE1:}} Datos de entrada de alta de contacto \\
\hspace*{1cm} Nombre: Cadena de caracteres (20) \\
\hspace*{1cm} Apellidos: Cadena de caracteres (40) \\
\hspace*{1cm} Teléfono: Cadena de caracteres (20) \\[12pt]

\textcolor{azul}{\textbf{RDW1:}} Datos almacenados de contacto. \\
\hspace*{1cm} Los mismos datos que RDE1.

\part{Práctica 1}

\chapter*{Datos Iniciales}
\addcontentsline{toc}{chapter}{Datos Iniciales}

\begin{itemize}
  \item Nombre del sistema: Red Social.
  \item Nombre de los subsistemas y personas responsables de cada subsistema:
    \begin{itemize}
      \item Publicaciones: Javier Niño Sánchez.
      \item Tendencias: Jesús Rodríguez González.
      \item Publicidad: Ismael Sallami Moreno.
      \item Usuarios: Fernando JoséGracia Choin.
      \item Mensajería Privada: Sergio Calvo González.
    \end{itemize}

    \item Descripción del Sistema EKIS (Red Social): 

    EKIS es un \textbf{Sistema de Información de Red Social} diseñado para permitir y gestionar la interacción y el intercambio de información y contenido entre sus usuarios. Su misión fundamental es proporcionar una plataforma robusta para la comunicación social, la creación de perfiles de usuario y la difusión de contenido en diversos formatos. El sistema está intrínsecamente ligado a la gestión de datos, requiriendo un análisis exhaustivo de los agentes externos que activan las funcionalidades y los distintos tipos de datos que intervienen en las tareas del sistema. Estructuralmente, EKIS se divide en cinco subsistemas independientes pero interconectados: Publicaciones, Tendencias, Publicidad, Usuarios y Mensajería Privada. El sistema debe garantizar la integridad de los datos a través de la aplicación de restricciones semánticas específicas que describen cambios en el comportamiento de las funciones ante configuraciones de datos predefinidas. En resumen, EKIS es una plataforma integral que equilibra la necesidad de comunicación social y gestión de contenido con la infraestructura necesaria para, si fuera necesario, integrar sistemas como la publicidad y el análisis de la popularidad de temas de la actualidad. \\

    \textbf{Descripción de los Subsistemas}

    \begin{enumerate}
      \item \textbf{Subsistema de Publicaciones (Responsable: Javier Niño Sánchez)}

      El subsistema de Publicaciones tiene la responsabilidad central de administrar todo el contenido generado por los miembros de la red social. Su función principal es el área de la difusión de contenido, permitiendo a los usuarios la creación de nuevas entradas, las cuales pueden incluir nombres, descripciones, categorías e imágenes opcionales. Este subsistema se encarga de que los usuarios puedan modificar posteriormente el contenido que han creado, asegurando que solo el propietario de la publicación pueda hacerlo, y de gestionar las interacciones de la comunidad, específicamente la funcionalidad de “Me gusta”. Una responsabilidad clave es la presentación de la información, permitiendo a los usuarios visualizar las publicaciones de aquellos a quienes siguen. Además, el subsistema garantiza la trazabilidad de los datos, ya que la eliminación de una publicación no resulta en un borrado físico, sino en un marcado que la oculta de la vista pública.\\

      \item \textbf{Subsistema de Tendencias (Responsable: Jesús Rodríguez González)}

      El subsistema de Tendencias es esencialmente el módulo de análisis de popularidad de EKIS, encargado de identificar y monitorizar los temas más discutidos en la plataforma. Su responsabilidad principal es el seguimiento y el conteo de hashtags, activándose periódicamente para revisar si se han mencionado hashtags en las últimas publicaciones de la plataforma. Este subsistema mantiene un registro actualizado de la frecuencia de aparición de cada hashtag. Entre sus funciones críticas se encuentra la capacidad de mostrar a los usuarios los temas más candentes (los 10 hashtags con más apariciones). Adicionalmente, el subsistema otorga al administrador la potestad de clasificar y organizar estos temas populares, asignándoles categorías específicas, y permite la consulta ordenada de tendencias por categoría. Los administradores también podrán ''eliminar'' un tema popular reseteando su contador de apariciones a cero.\\

      \item \textbf{Subsistema de Publicidad (Responsable: Ismael Sallami Moreno)}

      Este subsistema opera como el centro de gestión de anuncios de EKIS, centrándose en la administración eficiente de los anuncios que se presentan a los usuarios. Las responsabilidades son ejecutadas principalmente por el administrador e incluyen el control del ciclo de vida de los anuncios. Esto implica la creación de nuevos anuncios, donde se definen parámetros cruciales como el identificador único, el contenido (título y cuerpo), el enlace de destino y el periodo de vigencia específico (Fechas de Inicio y Fin). Una de las funciones más importantes es la segmentación, que permite al administrador definir el público objetivo al asignar características específicas a cada anuncio. El subsistema también debe manejar la modificación de los detalles de un anuncio ya existente y la eliminación del mismo, con la restricción de que los anuncios activos (dentro de su periodo de vigencia) no pueden ser borrados. Finalmente, tiene la responsabilidad de proporcionar la información completa de un anuncio al sistema cuando sea necesario para su visualización o procesamiento.\\

      \item \textbf{Subsistema de Usuarios (Responsable: Fernando José Gracia Choin)}

      El subsistema de Usuarios es el encargado de la identidad digital y la estructura de las relaciones sociales dentro de EKIS. Su ámbito abarca desde el proceso inicial de alta de un nuevo miembro, registrando información como nombre de usuario, correo electrónico y contraseña, hasta la gestión continua del perfil, permitiendo la modificación de datos personales, imágenes y biografías. La eliminación de cuentas es otra responsabilidad clave, asegurando que todos los datos personales asociados sean borrados de forma segura. En cuanto a las interacciones sociales, el subsistema tiene la responsabilidad dual de gestionar las relaciones positivas (como añadir a otros usuarios como amigos) y las relaciones restrictivas (como el bloqueo de otros usuarios para impedir la interacción mutua). Este subsistema debe aplicar rigurosas restricciones semánticas para evitar auto-bloqueos, duplicidad de relaciones y garantizar que los bloqueos y amistades se mantengan coherentes con la existencia de las cuentas.\\

      \item \textbf{Subsistema de Mensajería Privada (Responsable: Sergio Calvo González)}

      El subsistema de Mensajería Privada se dedica a facilitar la comunicación bidireccional, punto a punto y de carácter confidencial entre los usuarios de EKIS. La principal restricción de este subsistema es que solo está diseñado para operar entre usuarios que han establecido una relación de amistad. Sus responsabilidades esenciales incluyen el envío de mensajes de texto, asegurando que los mensajes no estén vacíos y que el destinatario exista en la plataforma. También se encarga de la eliminación de mensajes, garantizando que el texto borrado se vuelva inaccesible tanto para el emisor como para el receptor. Además de la transmisión de mensajes, el subsistema es responsable de la organización y presentación de la comunicación, permitiendo a los usuarios listar todas sus conversaciones activas y visualizar el historial de mensajes con un contacto específico de forma ordenada. Finalmente, ofrece funcionalidad para que los usuarios puedan personalizar la visibilidad de sus chats mediante la opción de archivar o desarchivar conversaciones.\\
    \end{enumerate}

    % \item Listado de requisitos funcionales, requisitos de datos y restricciones semánticas:
\end{itemize}


\chapter{Publicaciones}
\section{Primer requisito funcional}
\textbf{\textcolor{rojo}{\underline{RF1:}}} Crear comunidad. \\[6pt]

\textbf{\textcolor{rojo}{Entrada:}} Agente externo: usuario. Acción: solicitar inserción. 
Requisito de datos de entrada \textcolor{azul}{RDE2.1}. \\[6pt]

\textbf{\textcolor{rojo}{BD:}} Requisito de datos de escritura \textcolor{azul}{RDW2.1}. \\[6pt]

\textbf{\textcolor{rojo}{Salida:}} Agente externo: usuario. Acción: confirmación resultado. 
Requisito de datos de salida: ninguno. \\[12pt]

\textcolor{azul}{\textbf{RDE2.1:}} Datos de entrada de alta de contacto \\
\hspace*{1cm} Nombre de comunidad: Cadena de caracteres (20) \\
\hspace*{1cm} Imagen: Cadena de caracteres (255) \\
\hspace*{1cm} Descripción: Cadena de caracteres (80) \\ 
\hspace*{1cm} Categoría: Cadena de caracteres (20) \\[12pt]

\textcolor{azul}{\textbf{RDW2.1:}} Datos almacenados de contacto. \\
\hspace*{1cm} ID de comunidad: Cadena de caracteres (40) \\
\hspace*{1cm} Fecha de creación: Cadena de caracteres (40) \\
\pagebreak

\section{Segundo requisito funcional}
\textbf{\textcolor{rojo}{\underline{RF1:}}} Modificar comunidad. \\[6pt]

\textbf{\textcolor{rojo}{Entrada:}} Agente externo: usuario. Acción: solicitar inserción. 
Requisito de datos de entrada \textcolor{azul}{RDE1}. \\[6pt]

\textbf{\textcolor{rojo}{BD:}} Requisito de datos de escritura \textcolor{azul}{RDW1}. \\[6pt]

\textbf{\textcolor{rojo}{Salida:}} Agente externo: usuario. Acción: confirmación resultado. 
Requisito de datos de salida: ninguno. \\[12pt]

\textcolor{azul}{\textbf{RDE1:}} Datos de entrada de alta de contacto \\
\hspace*{1cm} Nombre: Cadena de caracteres (20) \\
\hspace*{1cm} Apellidos: Cadena de caracteres (40) \\
\hspace*{1cm} Teléfono: Cadena de caracteres (20) \\[12pt]

\textcolor{azul}{\textbf{RDW1:}} Datos almacenados de contacto. \\
\hspace*{1cm} Los mismos datos que RDE1.

\section{Tercer requisito funcional}
\textbf{\textcolor{rojo}{\underline{RF1:}}} Listar comunidades. \\[6pt]

\textbf{\textcolor{rojo}{Entrada:}} Agente externo: usuario. Acción: solicitar inserción. 
Requisito de datos de entrada \textcolor{azul}{RDE1}. \\[6pt]

\textbf{\textcolor{rojo}{BD:}} Requisito de datos de escritura \textcolor{azul}{RDW1}. \\[6pt]

\textbf{\textcolor{rojo}{Salida:}} Agente externo: usuario. Acción: confirmación resultado. 
Requisito de datos de salida: ninguno. \\[12pt]

\textcolor{azul}{\textbf{RDE1:}} Datos de entrada de alta de contacto \\
\hspace*{1cm} Nombre: Cadena de caracteres (20) \\
\hspace*{1cm} Apellidos: Cadena de caracteres (40) \\
\hspace*{1cm} Teléfono: Cadena de caracteres (20) \\[12pt]

\textcolor{azul}{\textbf{RDW1:}} Datos almacenados de contacto. \\
\hspace*{1cm} Los mismos datos que RDE1.

\section{Cuarto requisito funcional}
\textbf{\textcolor{rojo}{\underline{RF1:}}} Añadir publicación. \\[6pt]

\textbf{\textcolor{rojo}{Entrada:}} Agente externo: usuario. Acción: solicitar inserción. 
Requisito de datos de entrada \textcolor{azul}{RDE1}. \\[6pt]

\textbf{\textcolor{rojo}{BD:}} Requisito de datos de escritura \textcolor{azul}{RDW1}. \\[6pt]

\textbf{\textcolor{rojo}{Salida:}} Agente externo: usuario. Acción: confirmación resultado. 
Requisito de datos de salida: ninguno. \\[12pt]

\textcolor{azul}{\textbf{RDE1:}} Datos de entrada de alta de contacto \\
\hspace*{1cm} Nombre: Cadena de caracteres (20) \\
\hspace*{1cm} Apellidos: Cadena de caracteres (40) \\
\hspace*{1cm} Teléfono: Cadena de caracteres (20) \\[12pt]

\textcolor{azul}{\textbf{RDW1:}} Datos almacenados de contacto. \\
\hspace*{1cm} Los mismos datos que RDE1.

\section{Quinto requisito funcional}
\textbf{\textcolor{rojo}{\underline{RF1:}}} Expulsar usuario. \\[6pt]

\textbf{\textcolor{rojo}{Entrada:}} Agente externo: usuario. Acción: solicitar inserción. 
Requisito de datos de entrada \textcolor{azul}{RDE1}. \\[6pt]

\textbf{\textcolor{rojo}{BD:}} Requisito de datos de escritura \textcolor{azul}{RDW1}. \\[6pt]

\textbf{\textcolor{rojo}{Salida:}} Agente externo: usuario. Acción: confirmación resultado. 
Requisito de datos de salida: ninguno. \\[12pt]

\textcolor{azul}{\textbf{RDE1:}} Datos de entrada de alta de contacto \\
\hspace*{1cm} Nombre: Cadena de caracteres (20) \\
\hspace*{1cm} Apellidos: Cadena de caracteres (40) \\
\hspace*{1cm} Teléfono: Cadena de caracteres (20) \\[12pt]

\textcolor{azul}{\textbf{RDW1:}} Datos almacenados de contacto. \\
\hspace*{1cm} Los mismos datos que RDE1.


\chapter{Tendencias}
% Listado de requisitos funcionales, requisitos de datos y restricciones semánticas.

% Funcionales: 
% \begin{itemize}
%     \item Crear o mencionar hashtag
%     \item Listar tendencias
%     \item Asignar una categoría a un hashtag
%     \item Mostrar categoría ordenada
%     \item Eliminar tendencia
% \end{itemize}

%======================= %
% REQUISITOS FUNCIONALES %
% ======================= %


\section{Crear o mencionar hashtag}
\textbf{\textcolor{rojo}{\underline{RF2.1:}}} Crear o mencionar hashtag. \\[6pt]

\textbf{\textcolor{rojo}{\underline{Descripción:}}} Cada vez que haya una publicación con una cadena que empiece por "\#", almacenarlo de cara a poder ver los temas que son tendencia. En caso de ya existir, se actualiza el contador de veces que ha aparecido ese hashtag; si no, se crea y se pone el contador a uno.
 \\[6pt]

\textbf{\textcolor{rojo}{Entrada:}} Agente externo: usuario. Acción: realizar una publicación que contenga una cadena que empiece por "\#". 
Requisito de datos de entrada \textcolor{azul}{RDE2.1}. \\[6pt]

\textbf{\textcolor{rojo}{BD:}} Requisito de datos de escritura \textcolor{azul}{RDW2.1} y lectura \textcolor{azul}{RDR2.1.} \\[6pt]

\textbf{\textcolor{rojo}{Salida:}} Ninguna \\[6pt]

\textcolor{azul}{\textbf{RDE2.1:}} Datos de entrada de crear o mencionar hashtag \\
\hspace*{1cm} Hashtag: Cadena de caracteres (32) \\[12pt]

\textcolor{azul}{\textbf{RDR2.1:}} Datos para comprobar si ya existía un hashtag. \\
\hspace*{1cm} Los mismos que RDE2.1\\

\textcolor{azul}{\textbf{RDW2.1:}} Datos almacenados/actualizados de un hashtag. \\
\hspace*{1cm} Los mismos que RDE2.1 además de el contador de veces que ha aparecido ese hashtag.\\
\hspace*{1cm} Contador: Entero \\[12pt]

\textbf{\textcolor{verde}{Restricciones Semánticas:}} \\[6pt]
\textbf{\textcolor{verde}{\underline{RS2.1:}}} El Hastag debe empezar por "\#". \textbf{RF:} RF2.1. \textbf{RD(s):} RDW2.1. \\
\hspace*{1cm} \textit{Descripción:} Si el hashtag no había sido creado y no empieza por "\#", no se inserta el nuevo hashtag y se devuelve un error. \\[6pt]

\textbf{\textcolor{verde}{\underline{RS2.2:}}} El Hashtag puede ya existir o ser uno nuevo. \textbf{RF:} RF2.1. \textbf{RD(s):} RDW2.1 y RDR2.1. \\
\hspace*{1cm} \textit{Descripción:} El hashtag introducido puede que no existiera aún en la base de datos o que si existiera. En caso de ya existir, se actualiza el contador de veces que ha aparecido ese hashtag sumándole uno; si no, se crea y se pone el contador a uno.\\[6pt]
\pagebreak

\section{Listar tendencias}
\textbf{\textcolor{rojo}{\underline{RF2.2:}}} Listar tendencias. \\[6pt]

\textbf{\textcolor{rojo}{\underline{Descripción:}}} Mostrar los 10 hashtags con mayores apariciones.
 \\[6pt]

\textbf{\textcolor{rojo}{Entrada:}} Agente externo: usuario. Acción: solicitar listado de tendencias. 
Requisito de datos de entrada: ninguno. \\[6pt]

\textbf{\textcolor{rojo}{BD:}} Requisito de datos de lectura \textcolor{azul}{RDR2.2}\\[6pt]

\textbf{\textcolor{rojo}{Salida:}} Agente externo: usuario. Acción: confirmación del resultado.
Requisitos de datos de salida: \textcolor{azul}{RDS2.2}. \\[6pt] 

\textcolor{azul}{\textbf{RDR2.2:}} Datos para ver cuáles son los hashtags con más apariciones. \\
\hspace*{1cm} Los mismos que RDW2.1\\

\textcolor{azul}{\textbf{RDS2.1:}} Datos para listar tendencias. \\
\hspace*{1cm} Los mismos que RDW2.1\\

\pagebreak

\section{Asignar una categoría a un hashtag}
\textbf{\textcolor{rojo}{\underline{RF2.3:}}} Asignar una categoría a un hashtag. \\[6pt]

\textbf{\textcolor{rojo}{\underline{Descripción:}}} Asignar una categoría a un hashtag ya existente para poder clasificarlos.\\[6pt]

\textbf{\textcolor{rojo}{Entrada:}} Agente externo: administrador. Acción: solicitar actualizar la categoría de un hashtag. 
Requisito de datos de entrada \textcolor{azul}{RDE2.3} \\[6pt]

\textbf{\textcolor{rojo}{BD:}} Requisito de datos de lectura \textcolor{azul}{RDR2.3} y de escritura \textcolor{azul}{RDW2.3}\\[6pt]

\textbf{\textcolor{rojo}{Salida:}} Ninguna\\[6pt] 

\textcolor{azul}{\textbf{RDE2.3:}} El hashtag y la categoría que se le quiere asignar. \\
\hspace*{1cm} Hashtag: Cadena de caracteres (32)\\[12pt]
\hspace*{1cm} Categoría: Cadena de caracteres (32)\\[12pt]

\textcolor{azul}{\textbf{RDR2.3:}} Hashtag que hay que comprobar si ya existe en el sistema. \\
\hspace*{1cm} El mismo que \textcolor{azul}{RDR2.1}\\

\textcolor{azul}{\textbf{RDW2.3:}} Categoría a asignar. \\
\hspace*{1cm} Categoría: Cadena de caracteres (32)\\[12pt]

\textbf{\textcolor{verde}{Restricciones Semánticas:}} \\[6pt]
\textbf{\textcolor{verde}{\underline{RS2.3:}}} El hashtag debe existir previamente en el sistema. \textbf{RF:} RF2.3. \textbf{RD(s):} RDW2.3 y RDR2.3. \\
\hspace*{1cm} \textit{Descripción:} Si el hashtag no se encuentra almacenado en el sistema, no se crea ni se le asigna la categoría, sino que se devuelve un error. \\[6pt]

\pagebreak

\section{Mostrar categoría ordenada}
\textbf{\textcolor{rojo}{\underline{RF2.4:}}} Mostrar categoría ordenada. \\[6pt]

\textbf{\textcolor{rojo}{\underline{Descripción:}}} Listar las tendencias de una determinada categoría ordenadas de mayor a menos número de apariciones.\\[6pt]

\textbf{\textcolor{rojo}{Entrada:}} Agente externo: usuario. Acción: solicitar mostrar una categoría. 
Requisito de datos de entrada \textcolor{azul}{RDE2.4} \\[6pt]

\textbf{\textcolor{rojo}{BD:}} Requisito de datos de lectura \textcolor{azul}{RDR2.4}\\[6pt]

\textbf{\textcolor{rojo}{Salida:}} Agente externo: usuario. Acción: confirmación del resultado.
Requisitos de datos de salida: \textcolor{azul}{RDS2.4}. \\[6pt] 

\textcolor{azul}{\textbf{RDE2.4:}} La categoría que se quiere mostrar. \\
\hspace*{1cm} El mismo que \textcolor{azul}{RDW2.3}\\

\textcolor{azul}{\textbf{RDR2.4:}} Datos necesarios para ver si un hashtag pertenece a la categoría dada y ordenarlos por su número de apariciones. \\
\hspace*{1cm} Los mismos que \textcolor{azul}{RDE2.3} +  \textcolor{azul}{RDE2.4}\\

\textcolor{azul}{\textbf{RDS2.4:}} Datos para listar la categoría. \\
\hspace*{1cm} Los mismos que  \textcolor{azul}{RDR2.4}\\

\textbf{\textcolor{verde}{Restricciones Semánticas:}} \\[6pt]
\textbf{\textcolor{verde}{\underline{RS2.4:}}} La categoría debe existir previamente en el sistema. \textbf{RF:} RF2.4. \textbf{RD(s):} RDR2.4. \\
\hspace*{1cm} \textit{Descripción:} Si la categoría no se encuentra almacenada en el sistema se devuelve un error. \\[6pt]

\pagebreak

\section{Eliminar tendencia}
\textbf{\textcolor{rojo}{\underline{RF2.5:}}} Eliminar tendencia. \\[6pt]

\textbf{\textcolor{rojo}{\underline{Descripción:}}} Hacer que un hashtag deje de ser tendencia poniendo su contador de apariciones a cero.\\[6pt]

\textbf{\textcolor{rojo}{Entrada:}} Agente externo: administrador. Acción: solicitar la eliminación de una tendencia. 
Requisito de datos de entrada \textcolor{azul}{RDE2.5} \\[6pt]

\textbf{\textcolor{rojo}{BD:}} Requisito de datos de lectura \textcolor{azul}{RDR2.5} y de escritura \textcolor{azul}{RDW2.5}\\[6pt]

\textbf{\textcolor{rojo}{Salida:}} Ninguna

\textcolor{azul}{\textbf{RDE2.5:}} Hashtag que se quiere eliminar. \\
\hspace*{1cm} El mismo que \textcolor{azul}{RDE2.1}\\

\textcolor{azul}{\textbf{RDW2.5:}} Datos que hay que modificar para hacer que el hashtag deje de ser tendencia. \\
\hspace*{1cm} Contador: Entero\\[12pt]

\textcolor{azul}{\textbf{RDR2.5:}} Hashtag que se quiere eliminar. \\
\hspace*{1cm} El mismo que \textcolor{azul}{RDE2.1}\\

\textbf{\textcolor{verde}{Restricciones Semánticas:}} \\[6pt]
\textbf{\textcolor{verde}{\underline{RS2.5:}}} El hashtag debe existir previamente en el sistema. \textbf{RF:} RF2.5. \textbf{RD(s):} RDW2.5 y RDR2.5. \\
\hspace*{1cm} \textit{Descripción:} Si el hashtag no se encuentra almacenado en el sistema se devuelve un error. \\[6pt]

\pagebreak



\chapter{Publicidad}
% Listado de requisitos funcionales, requisitos de datos y restricciones semánticas.

% Funcionales: 
% \begin{itemize}
%     \item Añadir anuncio
%     \item Añadir característica a un anuncio del público objetivo
%     \item Eliminar anuncio
%     \item Modificar anuncio
%     \item Consultar anuncio
%     \item Redirigir a enlace externo
% \end{itemize}

%======================= %
% REQUISITOS FUNCIONALES %
% ======================= %


\section{Añadir anuncio}
\textbf{\textcolor{rojo}{\underline{RF1:}}} Añadir anuncio: Permite a un administrador crear un nuevo anuncio en el sistema proporcionando toda la información necesaria, como el título, cuerpo, enlace de destino y período de vigencia. \\[6pt]

\textbf{\textcolor{rojo}{Entrada:}} Agente externo: administrador. Acción: solicitar inserción de un nuevo anuncio. 
Requisito de datos de entrada \textcolor{azul}{RDE1}. \\[6pt]

\textbf{\textcolor{rojo}{BD:}} Requisito de datos de escritura \textcolor{azul}{RDW1}. \\[6pt]

\textbf{\textcolor{rojo}{Salida:}} Agente externo: administrador. Acción: confirmación del resultado de la inserción. 
Requisito de datos de salida: ninguno. \\[12pt]

\textcolor{azul}{\textbf{RDE1:}} Datos de entrada de alta de anuncio \\
\hspace*{1cm} Identificador Anuncio: Cadena de caracteres (10) \\
\hspace*{1cm} Título: Cadena de caracteres (100) \\
\hspace*{1cm} Cuerpo del Anuncio: Cadena de caracteres (500) \\
\hspace*{1cm} Enlace Destino: Cadena de caracteres (255) \\
\hspace*{1cm} Fecha de Inicio: Fecha \\
\hspace*{1cm} Fecha de Fin: Fecha \\[12pt]

\textcolor{azul}{\textbf{RDW1:}} Datos almacenados del anuncio. \\
\hspace*{1cm} Los mismos datos que RDE1. \\[12pt]

\textbf{\textcolor{verde}{Restricciones Semánticas:}} \\[6pt]
\textbf{\textcolor{verde}{\underline{RS1.1:}}} El Identificador del Anuncio debe ser único. \textbf{RF:} RF1. \textbf{RD(s):} RDW1. \\
\hspace*{1cm} \textit{Descripción:} Si ya existe un anuncio con el mismo identificador, no se inserta el nuevo anuncio y se devuelve un error. \\[6pt]

\textbf{\textcolor{verde}{\underline{RS1.2:}}} La fecha de fin debe ser posterior a la fecha de inicio. \textbf{RF:} RF1. \textbf{RD(s):} RDE1. \\
\hspace*{1cm} \textit{Descripción:} Si la fecha de fin no es posterior a la fecha de inicio, la inserción es rechazada y se devuelve un mensaje de error.

\section{Añadir característica a un anuncio del público objetivo}
\textbf{\textcolor{rojo}{\underline{RF2:}}} Añadir característica a un anuncio del público objetivo: Permite asignar una característica específica a un anuncio para definir su público objetivo. Por ejemplo, se puede asignar la característica ''Edad'' con el valor ''18-25'' para dirigir el anuncio a usuarios dentro de ese rango de edad. \\[6pt]

\textbf{\textcolor{rojo}{Entrada:}} Agente externo: administrador. Acción: solicitar asignación de característica. 
Requisito de datos de entrada \textcolor{azul}{RDE2}. \\[6pt]

\textbf{\textcolor{rojo}{BD:}} Requisito de datos de escritura \textcolor{azul}{RDW2}. \\[6pt]

\textbf{\textcolor{rojo}{Salida:}} Agente externo: administrador. Acción: confirmación de la asignación. 
Requisito de datos de salida: ninguno. \\[12pt]

\textcolor{azul}{\textbf{RDE2:}} Datos para añadir característica a un anuncio \\
\hspace*{1cm} Identificador Anuncio: Cadena de caracteres (10) \\
\hspace*{1cm} Nombre Característica: Cadena de caracteres (50) \\
\hspace*{1cm} Valor Característica: Cadena de caracteres (100) \\[12pt]

\textcolor{azul}{\textbf{RDW2:}} Datos almacenados de la característica del público objetivo para un anuncio. \\
\hspace*{1cm} Los mismos datos que RDE2. \\[12pt]


\textbf{\textcolor{verde}{Restricciones Semánticas:}} \\[6pt]
\textbf{\textcolor{verde}{\underline{RS2.1:}}} El anuncio debe existir. \textbf{RF:} RF2. \textbf{RD(s):} RDE2. \\
\hspace*{1cm} \textit{Descripción:} Si el `Identificador Anuncio` no corresponde a un anuncio existente en la base de datos, no se realiza la inserción y se devuelve un aviso.

\section{Eliminar anuncio}
\textbf{\textcolor{rojo}{\underline{RF3:}}} Eliminar anuncio: Permite a un administrador eliminar un anuncio existente del sistema utilizando su identificador único. \\[6pt]

\textbf{\textcolor{rojo}{Entrada:}} Agente externo: administrador. Acción: solicitar borrado de un anuncio. 
Requisito de datos de entrada \textcolor{azul}{RDE3}. \\[6pt]

\textbf{\textcolor{rojo}{BD:}} Requisito de datos de escritura \textcolor{azul}{RDW3}. \\[6pt]

\textbf{\textcolor{rojo}{Salida:}} Agente externo: administrador. Acción: confirmación del resultado del borrado. 
Requisito de datos de salida: ninguno. \\[12pt]

\textcolor{azul}{\textbf{RDE3:}} Datos de entrada de baja de anuncio \\
\hspace*{1cm} Identificador Anuncio: Cadena de caracteres (10) \\[12pt]

\textcolor{azul}{\textbf{RDW3:}} Identificador del anuncio a eliminar. \\
\hspace*{1cm} Los mismos datos que RDE3. \\[12pt] 

\textbf{\textcolor{verde}{Restricciones Semánticas:}} \\[6pt]
\textbf{\textcolor{verde}{\underline{RS3.1:}}} No se puede eliminar un anuncio activo. \textbf{RF:} RF3. \textbf{RD(s):} RDE3. \\
\hspace*{1cm} \textit{Descripción:} Si la fecha actual está dentro del período de vigencia del anuncio (entre Fecha de Inicio y Fecha de Fin), no se permite el borrado y se devuelve un error.

\section{Modificar anuncio}
\textbf{\textcolor{rojo}{\underline{RF4:}}} Modificar anuncio: Permite a un administrador actualizar la información de un anuncio existente, como el título, cuerpo, enlace de destino y período de vigencia. Los campos son opcionales, permitiendo modificar solo los necesarios. \\[6pt]

\textbf{\textcolor{rojo}{Entrada:}} Agente externo: administrador. Acción: solicitar modificación de un anuncio existente. 
Requisito de datos de entrada \textcolor{azul}{RDE4}. \\[6pt]

\textbf{\textcolor{rojo}{BD:}} Requisito de datos de escritura \textcolor{azul}{RDW4}. \\[6pt]

\textbf{\textcolor{rojo}{Salida:}} Agente externo: administrador. Acción: confirmación del resultado de la modificación. 
Requisito de datos de salida: ninguno. \\[12pt]

\textcolor{azul}{\textbf{RDE4:}} Datos de entrada para modificar un anuncio \\
\hspace*{1cm} Identificador Anuncio: Cadena de caracteres (10) \\
\hspace*{1cm} Título (opcional): Cadena de caracteres (100) \\
\hspace*{1cm} Cuerpo del Anuncio (opcional): Cadena de caracteres (500) \\
\hspace*{1cm} Enlace Destino (opcional): Cadena de caracteres (255) \\
\hspace*{1cm} Fecha de Inicio (opcional): Fecha \\
\hspace*{1cm} Fecha de Fin (opcional): Fecha \\[12pt]

\textcolor{azul}{\textbf{RDW4:}} Datos del anuncio a actualizar. \\
\hspace*{1cm} Los mismos datos que RDE4. \\[12pt]


\textbf{\textcolor{verde}{Restricciones Semánticas:}} \\[6pt]
\textbf{\textcolor{verde}{\underline{RS4.1:}}} La fecha de fin debe ser posterior a la fecha de inicio. \textbf{RF:} RF4. \textbf{RD(s):} RDE4. \\
\hspace*{1cm} \textit{Descripción:} Si se modifican las fechas y la nueva fecha de fin no es posterior a la nueva fecha de inicio, la modificación es rechazada y se devuelve un mensaje de error.

\section{Consultar anuncio}
\textbf{\textcolor{rojo}{\underline{RF5:}}} Consultar anuncio: Permite al sistema recuperar toda la información de un anuncio específico utilizando su identificador único, para mostrarlo o procesarlo según sea necesario.\\[6pt]

\textbf{\textcolor{rojo}{Entrada:}} Agente externo: sistema. Acción: solicitar datos de un anuncio por su identificador. 
Requisito de datos de entrada \textcolor{azul}{RDE5}. \\[6pt]

\textbf{\textcolor{rojo}{BD:}} Requisito de datos de lectura \textcolor{azul}{RDR5}. \\[6pt]

\textbf{\textcolor{rojo}{Salida:}} Agente externo: sistema. Acción: entregar datos del anuncio. 
Requisito de datos de salida: \textcolor{azul}{RDS5}. \\[12pt]

\textcolor{azul}{\textbf{RDE5:}} Datos de entrada para consultar un anuncio \\
\hspace*{1cm} Identificador Anuncio: Cadena de caracteres (10) \\[12pt]

\textcolor{azul}{\textbf{RDR5:}} Datos del anuncio a leer. \\
\hspace*{1cm} Todos los datos asociados al `Identificador Anuncio`, equivalentes a RDW1. \\[12pt]

\textcolor{azul}{\textbf{RDS5:}} Datos de salida del anuncio consultado. \\
\hspace*{1cm} Los mismos datos que RDR5. \\[12pt]
 
\textbf{\textcolor{verde}{Restricciones Semánticas:}} \\[6pt]
\textbf{\textcolor{verde}{\underline{RS5.1:}}} El anuncio debe existir. \textbf{RF:} RF5. \textbf{RD(s):} RDE5. \\
\hspace*{1cm} \textit{Descripción:} Si el `Identificador Anuncio` no existe, la consulta devuelve un resultado vacío o un mensaje de error.

% \section{Redirigir a enlace externo}
% \textbf{\textcolor{rojo}{\underline{RF6:}}} Redirigir a enlace externo. \\[6pt]

% \textbf{\textcolor{rojo}{Entrada:}} Agente externo: usuario final. Acción: hacer clic en un anuncio. 
% Requisito de datos de entrada \textcolor{azul}{RDE6}. \\[6pt]

% \textbf{\textcolor{rojo}{BD:}} Requisito de datos de lectura \textcolor{azul}{RDR6}. \\[6pt]

% \textbf{\textcolor{rojo}{Salida:}} Agente externo: navegador del usuario final. Acción: redirigir a una URL. 
% Requisito de datos de salida: \textcolor{azul}{RDS6}. \\[12pt]

% \textcolor{azul}{\textbf{RDE6:}} Datos de entrada para la redirección \\
% \hspace*{1cm} Identificador Anuncio: Cadena de caracteres (10) \\[12pt]

% \textcolor{azul}{\textbf{RDR6:}} Datos del enlace a obtener. \\
% \hspace*{1cm} Enlace Destino: Cadena de caracteres (255) \\[12pt]

% \textcolor{azul}{\textbf{RDS6:}} Datos de salida para la redirección. \\
% \hspace*{1cm} Los mismos datos que RDR6. \\[12pt]

% \textbf{Descripción del requisito:} Permite redirigir al usuario final a una URL externa cuando hace clic en un anuncio, utilizando el enlace de destino asociado al anuncio.

% \textbf{\textcolor{verde}{Restricciones Semánticas:}} \\[6pt]
% \textbf{\textcolor{verde}{\underline{RS6.1:}}} El enlace de destino debe ser una URL válida. \textbf{RF:} RF6. \textbf{RD(s):} RDR6. \\
% \hspace*{1cm} \textit{Descripción:} Si el `Enlace Destino` recuperado no es una URL con un formato válido, se debe gestionar el error, por ejemplo, no realizando la redirección.


\chapter{Usuarios}
\PassOptionsToPackage{table,dvipsnames}{xcolor}

\section{Primer requisito funcional}
\textbf{\textcolor{rojo}{\underline{RF1-U:}}} Crear usuario. 

\textbf{\textcolor{rojo}{\underline{Descripción:}}} Permite a un visitante crear un usuario en el sistema, asignándole un nombre de usuario y una contraseña.
Opcionalmente, puede añadir una imagen de perfil y/o una breve biografía. 

\textbf{\textcolor{rojo}{Entrada:}} Agente externo: visitante. Acción: solicitar inserción.
Requisito de datos de entrada \textcolor{azul}{RDE1.1}. 
\textbf{\textcolor{rojo}{BD:}} Requisito de datos de escritura \textcolor{azul}{RDW1.1}. 

\textbf{\textcolor{rojo}{Salida:}} Agente externo: visitante. Acción: confirmación resultado.
Requisito de datos de salida: ninguno.
\pagebreak

\textcolor{azul}{\textbf{RDE1.1:}} Datos de entrada de crear usuario. 
\hspace*{1cm} Nombre de usuario: Cadena de caracteres (20)
\hspace*{1cm} Correo electrónico: Cadena de caracteres (50)
\hspace*{1cm} Contraseña: Cadena de caracteres (20)
\hspace*{1cm} Imagen de perfil: Cadena de caracteres (255)
\hspace*{1cm} Biografía: Cadena de caracteres (100) 

\textcolor{azul}{\textbf{RDW1.1:}} Datos almacenados del usuario. 
\hspace*{1cm} Los mismos que RDE1.1 
\hspace*{1cm} ID de usuario: Cadena de caracteres (40) 
\hspace*{1cm} Fecha de creación: Cadena de caracteres (40) 
\newpage

\chapter{Mensajería Privada}
\newcounter{RDW} \setcounter{RDW}{1}
\newcounter{RDR} \setcounter{RDR}{1}
\newcounter{RDE} \setcounter{RDE}{1}
\newcounter{RDS} \setcounter{RDS}{1}
\newcounter{RF} \setcounter{RF}{1}
% Descripción del subsistema
\section*{Descripción}
El subsistema de Mensajería Privada se encarga de la comunicación privada entre usuarios, consistiendo principalmente en el envío, la eliminación y la visualización de mensajes que se dan entre dos usuarios del sistema, limitandose sólo a aquellos que son amigos. Se incluye la posibilidad de archivar chats o destacar mensajes para personalizar la visualización de los mensajes por parte del usuario. Parte desarrollada por \textbf{Sergio Calvo González}.
\newpage

% ======================= %
%     ENVIAR MENSAJE      %
% ======================= %
\section{Enviar mensaje}
\textbf{\textcolor{rojo}{\underline{RF5.\theRF:}}}. Enviar mensaje: Permite que un usuario envíe un mensaje de texto a otro usuario (deben ser amigos)\\[6pt]
\stepcounter{RF}

\textbf{\textcolor{rojo}{Entrada:}} Agente externo: usuario. Acción: solicitar envío de un mensaje. 
Requisito de datos de entrada \textcolor{azul}{RDE5.\theRDE}. \\[6pt]

\textbf{\textcolor{rojo}{BD:}} Requisito de datos de escritura \textcolor{azul}{RDW5.\theRDW}. \\[6pt]

\textbf{\textcolor{rojo}{Salida:}} Agente externo: usuario. Acción: confirmación resultado. 
Requisito de datos de salida: ninguno. \\[12pt]

\textcolor{azul}{\textbf{RDE5.\theRDE:}} Mensaje y usuario al que va dirigido \\
\hspace*{1cm} Usuario: Cadena de caracteres (20) \\
\hspace*{1cm} Mensaje: Cadena de caracteres (50) \\[12pt]
\stepcounter{RDE}

\textcolor{azul}{\textbf{RDW5.\theRDW:}} Datos del envío de mensaje. \\
\hspace*{1cm} Usuario Emisor: Cadena de caracteres (20) \\
\hspace*{1cm} Usuario Receptor: Cadena de caracteres (20) \\
\hspace*{1cm} Mensaje: Cadena de caracteres (50) \\[12pt]
\stepcounter{RDW}
% ======================= %
%    ELIMINAR MENSAJE     %
% ======================= %
\section{Eliminar mensaje}
\textbf{\textcolor{rojo}{\underline{RF\theRF:}}}. Eliminar mensaje: Permite que un usuario elimine un mensaje de texto que ha enviado, de forma que ni dicho usuario ni el receptor pueden volver a acceder al mensaje\\[6pt]
\stepcounter{RF}

\textbf{\textcolor{rojo}{Entrada:}} Agente externo: usuario. Acción: solicitar la eliminación de un mensaje. 
Requisito de datos de entrada: ninguno. \\[6pt]

\textbf{\textcolor{rojo}{BD:}} Requisito de datos de lectura \textcolor{azul}{RDR6.1}. \\[6pt]

\textbf{\textcolor{rojo}{Salida:}} Agente externo: usuario. Acción: confirmación resultado. 
Requisito de datos de salida: ninguno. \\[12pt]

\textcolor{azul}{\textbf{RDE6.2:}} Datos del mensaje a borrar\\
\hspace*{1cm} Usuario: Cadena de caracteres (20) \\
\hspace*{1cm} Mensaje: Cadena de caracteres (50) \\[12pt]

\textcolor{azul}{\textbf{RDR6.1:}} Lista de mensajes. \\
\hspace*{1cm} Usuario Emisor: Cadena de caracteres (20) \\
\hspace*{1cm} Usuario Receptor: Cadena de caracteres (20) \\
\hspace*{1cm} Mensaje: Cadena de caracteres (50) \\[12pt]

% ======================= %
%  LISTAR CONVERSACIONES  %
% ======================= %
\section{Listar conversaciones}
\textbf{\textcolor{rojo}{\underline{RF\theRF:}}}. Listar conversaciones: Permite que un usuario vea los conversaciones que tiene disponibles\\[6pt]
\stepcounter{RF}
\textbf{\textcolor{rojo}{Entrada:}} Agente externo: usuario. Acción: solicitar la visualización de sus conversaciones. 
Requisito de datos de entrada: ninguno\\[6pt]

\textbf{\textcolor{rojo}{BD:}} Requisito de datos de lectura \textcolor{azul}{RDR5.\theRDR}. \\[6pt]

\textbf{\textcolor{rojo}{Salida:}} Agente externo: usuario. Acción: confirmación resultado. 
Requisito de datos de salida: \textcolor{azul}{RDS5.\theRDS} \\[12pt]

\textcolor{azul}{\textbf{RDR5.\theRDR:}} Lista de usuarios con los que dispone de una conversación, con el formato:\\
\hspace*{1cm} Usuario: Cadena de caracteres (20) \\[12pt]
\stepcounter{RDR}

\textcolor{azul}{\textbf{RDS5.\theRDS:}} Lista de usuarios con los que dispone de una conversación. \\
\hspace*{1cm} Los mismos datos que (20) \\[12pt]
\stepcounter{RDS}
% ======================= %
%     VISUALIZAR CHAT     %
% ======================= %
\section{Visualizar chat}
\textbf{\textcolor{rojo}{\underline{RF\theRF:}}}. Visualizar chat: Permite que un usuario vea los mensajes entre otro usuario y él mismo\\[6pt]
\stepcounter{RF}
\textbf{\textcolor{rojo}{Entrada:}} Agente externo: usuario. Acción: solicitar la visualización de un chat. 
Requisito de datos de entrada \textcolor{azul}{RDE5.\theRDE}. \\[6pt]

\textbf{\textcolor{rojo}{BD:}} Requisito de datos de lectura \textcolor{azul}{RDR5.\theRDR}. \\[6pt]

\textbf{\textcolor{rojo}{Salida:}} Agente externo: usuario. Acción: confirmación resultado. 
Requisito de datos de salida: \textcolor{azul}{RDS6.2} \\[12pt]

\textcolor{azul}{\textbf{RDE5.\theRDE:}} Datos del usuario del chat que se desea ver\\
\hspace*{1cm} Usuario: Cadena de caracteres (20) \\[12pt]
\stepcounter{RDE}

\textcolor{azul}{\textbf{RDR5.\theRDR:}} Lista de mensajes entre ambos usuarios, ordenados temporalmente, con el siguiente formato:\\
\hspace*{1cm} Usuario emisor: Cadena de caracteres (20) \\
\hspace*{1cm} Mensaje: Cadena de caracteres (50) \\[12pt]


\textcolor{azul}{\textbf{RDS5.\theRDS:}} Lista de mensajes entre ambos usuarios.\\
\hspace*{1cm} Los mismos datos que en RDR5.\theRDR\\[12pt]
\stepcounter{RDR}
\stepcounter{RDS}
% ======================= %
%      ARCHIVAR CHAT      %
% ======================= %
\section{Des/archivar chat}
\textbf{\textcolor{rojo}{\underline{RF\theRF:}}}. Des/archivar chat: Permite que un usuario seleccione un chat del que quiere o no visualizar los mensajes, activando o no su visualización\\[6pt]
\stepcounter{RF}
\textbf{\textcolor{rojo}{Entrada:}} Agente externo: usuario. Acción: solicitar la des/archivación de un chat. 
Requisito de datos de entrada \textcolor{azul}{RDE5.\theRDE}. \\[6pt]

\textbf{\textcolor{rojo}{BD:}} Requisito de datos de escritura \textcolor{azul}{RDW5.\theRDW}. \\[6pt]

\textbf{\textcolor{rojo}{Salida:}} Agente externo: usuario. Acción: confirmación resultado. 
Requisito de datos de salida: ninguno \\[12pt]

\textcolor{azul}{\textbf{RDE5.\theRDE:}} Usuario que deseamos des/archivar   \\
\hspace*{1cm} Usuario: Cadena de caracteres (20) \\[12pt]

\textcolor{azul}{\textbf{RDW5.\theRDW:}} Datos del usuario que se desea archivar/desarchivar. \\
\hspace*{1cm} Usuario a archivar: Cadena de caracteres (20) \\
\hspace*{1cm} Usuario que archiva: Cadena de caracteres (20) \\
\hspace*{1cm} Usuario a archivar: Cadena de caracteres (20) \\[12pt]

\stepcounter{RDE}
\stepcounter{RDW}
% ======================== %
% RESTRICCIONES SEMÁNTICAS %
% ======================== %
\newcounter{RS} \setcounter{RS}{1}
\textbf{\textcolor{verde}{Restricciones Semánticas:}} \\[6pt]
\textbf{\textcolor{verde}{\underline{RS5.\theRS:}}} Los mensajes sólo pueden enviarse entre usuarios que son amigos. \\
\textbf{RF:} RF. \textbf{RD(s):} . \\
\hspace*{1cm} \textit{Descripción:} Si se desea enviar un mensaje a un usuario del que no se es amigo, se devolverá un error indicando que se necesita ser amigos para poder enviar un mensaje. \\[6pt]
\stepcounter{RS}

\textbf{\textcolor{verde}{\underline{RS5.\theRS:}}} Los mensajes no se pueden enviar al propio usuario. \\
\textbf{RF:} RF. \textbf{RD(s):} . \\
\hspace*{1cm} \textit{Descripción:} Si se desea enviar un mensaje al propio usuario, se devolverá un error indicando que no se puede enviar un mensaje a uno mismo. \\[6pt]
\stepcounter{RS}

\textbf{\textcolor{verde}{\underline{RS5.\theRS:}}} Para el envío de mensajes, el mensaje debe tener al menos un carácter. \\
\textbf{RF:} RF. \textbf{RD(s):} . \\
\hspace*{1cm} \textit{Descripción:} Si se desea enviar un mensaje, dicho mensaje no puede estar vacío. \\ [6pt]
\stepcounter{RS}

\textbf{\textcolor{verde}{\underline{RS5.\theRS:}}} Para el envío de mensajes, el usuario receptor debe existir. \\
\textbf{RF:} RF. \textbf{RD(s):} . \\
\hspace*{1cm} \textit{Descripción:} Si se desea enviar un mensaje, el receptor debe existir. \\ [6pt]
\stepcounter{RS}




% chapter{parte de cada uno}
% \input{chapters/nombre/parte.tex}

 \part{Práctica 2}

 \chapter{Publicaciones}
 
En este apartado abordaremos la parte correspondiente a la práctica 2 de la asignatura.
Se puede observar en la figura \ref{fig:dfd_publi} el Diagrama de Flujo de Datos (DFD) 
de nivel 1 del subsistema de publicidad. En este diagrama se representan los diferentes 
procesos, almacenes de datos, agentes externos y flujos de datos que componen el subsistema.

\begin{figure}[h!]
    \centering
    % Ajusta el tamaño del diagrama al ancho del documento (opcional)
    \shorthandoff{<>}
    \resizebox{\textwidth}{!}{%
    \begin{tikzpicture}[node distance=5cm, auto]

        % ======================================================
        % === 1. AGENTES EXTERNOS ===============================
        % ======================================================
        % Los "agentes" representan entidades externas al sistema
        % Ejemplo: usuarios, sistemas externos, administradores, etc.

        \node[agente, minimum width=4cm, minimum height=1.2cm] (Agente1) {Usuario};
        %\node[agente, right of=Agente1, node distance=14cm] (Agente2) {Sistema Externo};

        % ======================================================
        % === 2. ALMACÉN DE DATOS ===============================
        % ======================================================
        % Representan bases de datos o repositorios.
        % Suelen ser cilindros o rectángulos especiales.

        \node[almacen, below of=Agente1, node distance=10cm, xshift=7cm, minimum width=6cm] (BD1) {BD EKIS};

        % ======================================================
        % === 3. PROCESOS (FUNCIONES O RFs) =====================
        % ======================================================
        % Cada proceso es una función, requisito funcional (RFx.x) o tarea.
        % Se ubican entre los agentes y los datos.

        \node[proceso, below of=Agente1, node distance=3cm, xshift=-3cm, minimum width=4.5cm] (P1) {RF1.1: Crear publicación};
        \node[proceso, below of=P1, node distance=4cm, minimum width=4.5cm] (P2) {RF1.2: Modificar publicación};
        \node[proceso, right of=P1, node distance=8cm, minimum width=4.5cm] (P4) {RF1.4: Dar me gusta a una publicación};
        \node[proceso, right of=Agente1, node distance=6cm, xshift=3cm, minimum width=4.5cm] (P5) {RF1.5: Eliminar publicación};
        \node[proceso, below of=P4, node distance=2cm, minimum width=4.5cm] (P3) {RF1.3: Listar publicaciones};

        % ======================================================
        % === 4. FLUJOS DE DATOS ================================
        % ======================================================
        % Usa \draw[flujo] para representar conexiones.
        % Puedes añadir etiquetas con node[pos=..., ...]{Texto}.
        % Ejemplos:
        %   RDE = Requisito de entrada (data entrada)
        %   RDS = Requisito de salida (data salida)
        %   RDW = Escritura a BD
        %   RDR = Lectura desde BD

        % ------------------------------------------------------
        % Ejemplo de flujo entre Agente y Proceso (RDE)
        % ------------------------------------------------------
        \draw[flujo, ->] (Agente1.south west) ++(0.2,0) |- ++(0,-1) -| node[pos=0.7, right] {RDE 1.1} (P1.north);
        \draw[flujo, ->] (Agente1.west) |- ++(0,-1) -| node[pos=0.53] {RDE 1.2} (P2.north west);
        %RF1.3 no tiene RDE
        \draw[flujo, ->] (Agente1.south east) -| node[pos=0.6, left] {RDE 1.3} (P4.north);
        \draw[flujo, ->] (Agente1.east) -- node[pos=0.4, above] {RDE 1.4} (P5.west);
        
        % Proceso a Agente (RDS)
        %RF1.1 no tiene RDS
        \draw[flujo, ->] (P2.west) -- ++(-0.8,0) -| ++(-0.8,7) -- node[pos=0.4, above] {RDS 1.1} (Agente1.west);
        \draw[flujo, ->] (P3.west) -- node[pos=0.55, below] {RDS 1.2} ++(-1.5,0) |- ++(-0.97,2.5) -- (Agente1.south);
        \coordinate (elbow) at ($(P4.north west)+(0,1.875)$); 
        \draw[flujo, ->] (P4.north west) -- (elbow) node[pos=0.2, right] {RDS 1.3};

        % ------------------------------------------------------
        % Ejemplo de flujo entre Proceso y Almacén de Datos (RDW)
        % ------------------------------------------------------
        \draw[flujo, ->] (P1.east) -- ++(1.0,0) |- ++(1.0,-7) node[pos = 0.465] {RDW 1.1} -- (BD1.west);
        \draw[flujo, ->] (P2.south east) -- node[pos = 0.15, left] {RDW 1.2} ++(0, -3) -- (BD1.south west);
        %RF1.3 no genera ningún RDW
        \draw[flujo, ->] (P4.east) |- ++(-1.6, -3) -- node[pos = 0.1] {RDW 1.3} (BD1.north);
        \draw[flujo, ->] (P5.east) -- ++(0.5, 0) |- node[pos=0.68, below] {RDW 1.4} (BD1.east);
        %\draw[flujo, ->] (P5) -- (BD1); %ESTA LINEA

        % ------------------------------------------------------
        % Ejemplo de flujo desde Almacén a Proceso (RDR)
        % ------------------------------------------------------
        %RF1.1 no genera ningún RDR
        \draw[flujo, ->] (BD1.west) |- ++(-2,-1) -| node[pos=0.6, left] {RDR 1.1} (P2.south);
        \draw[flujo, ->] (BD1.north west) ++(0,-0.2) -- ++(0,2) node[pos=0.8, right] {RDR 1.2} -| (P3.south);
        \draw[flujo, ->] (BD1.north east) -- node[pos=0.2, left] {RDR 1.3} ++(0,6.49) -- (P4.east);

        % ------------------------------------------------------
        % Ejemplo de flujo entre Procesos
        % ------------------------------------------------------
        %\draw[flujo, ->] (P1.south) -- node[midway, right] {RDF interno} (P2.north);

        % ------------------------------------------------------
        % Ejemplo de flujo entre Proceso y Agente externo
        % ------------------------------------------------------
        %\draw[flujo, ->] (P4.east) -- ++(2,0) |- node[pos=0.3, right] {RDS 1.4} (Agente2.south);

    \end{tikzpicture}
    } % Fin del resizebox
    \caption{DFD1 — Subsistema de publicaciones}
    \label{fig:dfd_publi}
\end{figure}


 \chapter{Tendencias}
 
En este apartado abordaremos la parte correspondiente a la práctica 2 de la asignatura.
Se puede observar en la siguiente figura el Diagrama de Flujo de Datos (DFD) 
de nivel 1 del subsistema de tendencias. En este diagrama se representan los diferentes 
procesos, almacenes de datos, agentes externos y flujos de datos que componen el subsistema.

\begin{figure}[h!]
    \centering
    \shorthandoff{<>}
    \resizebox{\textwidth}{!}{%
    \begin{tikzpicture}[node distance=1.8cm and 4cm, auto]

        % === Agentes ===
        \node[agente, minimum width=4cm] (Usuario) {Usuario};
        \node[agente, below=of Usuario, yshift=-11cm] (Admin) {Admin};

        % === Almacén de datos ===
        \node[almacen, right=of Usuario, xshift=11.5cm, minimum width=5.5cm] (Hashtag) {Hashtag};

        % === Procesos ===
        \node[proceso, right=of Usuario, xshift=3cm] (RF21) {RF2.1: Crear o mencionar hashtag};
        \node[proceso, below=of RF21, yshift=-0.5cm] (RF22) {RF2.2: Listar tendencias};
        \node[proceso, below=of RF22, yshift=-0.5cm] (RF24) {RF2.4: Mostrar categoría ordenada};
        \node[proceso, below=of RF24, yshift=-0.5cm] (RF23) {RF2.3: Asignar categoría a un hashtag};
        \node[proceso, below=of RF23, yshift=-0.5cm] (RF25) {RF2.5: Eliminar tendencia};

        % === Flujo de datos ===

        % Usuario → RF2.1
        \draw[flujo, ->] (Usuario.east) -- node[above] {RDE 2.1} (RF21.west);

        % RF2.1 → Hashtag
        \draw[flujo, ->] (RF21.east) -- node[above] {RDW 2.1} (Hashtag.west);

        % Usuario → RF2.2
        \draw[flujo, ->] (Usuario.south) ++(1.2,0) |- node[pos=0.8, above left] {Consulta tendencias} (RF22.west);

        % RF2.2 → Usuario
        \draw[flujo, ->] (RF22.north) -- node[left] {RDS 2.2} (Usuario.south east);

        % Hashtag → RF2.2
        \draw[flujo, ->] (Hashtag.west) -- node[left] {RDR 2.2} (RF22.east);

        % Admin → RF2.3
        \draw[flujo, ->] (Admin.east) -- node[left] {RDE 2.3} (RF23.west);

        % RF2.3 → Hashtag
        \draw[flujo, ->] (RF23.east) -- node[right] {RDW 2.3} (Hashtag.south);

        % Hashtag → RF2.3
        \draw[flujo, ->] (Hashtag.south) ++(1,0) |- node[above left] {RDR 2.3} (RF23.south east);

        % Usuario → RF2.4
        \draw[flujo, ->] (Usuario.south) ++(0,-1.2) |- node[pos=0.8, above left] {RDE 2.4} (RF24.west);

        % Hashtag → RF2.4
        \draw[flujo, ->] (Hashtag.south) -- node[left] {RDR 2.4} (RF24.east);

        % Admin → RF2.5
        \draw[flujo, ->] (Admin.east) -- node[above] {RDE 2.5} (RF25.west);

        % RF2.5 → Hashtag
        \draw[flujo, ->] (RF25.east) |- ++(1.5,0) -| node[pos=0.8, above left] {RDW 2.5} (Hashtag.south east);


        % RF2.4 → Usuario
        \draw[flujo, ->] (RF24.west) |- ++(-2,-1) -| node[pos=0.6, left] {RDS 2.4} (Usuario.south west);

    \end{tikzpicture}
    }
    \caption{DFD1 — Subsistema de tendencias}
    \label{fig:dfd_tendencias}
\end{figure}


 \chapter{Publicidad}
 \chapter{Publicidad}

En este apartado abordaremos la parte correspondiente a la práctica 2 de la asginatura.

Se puede observar en la figura \ref{fig:dfd1_pub} el Diagrama de Flujo de Datos (DFD) de nivel 1 del subsistema de publicidad. En este diagrama se representan los diferentes procesos, almacenes de datos, agentes externos y flujos de datos que componen el subsistema.

\begin{figure}[h!]
    \centering
    \resizebox{\textwidth}{!}{%
    \begin{tikzpicture}[node distance=6cm, auto]

        % ==== 1. Agentes Externos ====
        \node[agente, minimum width=4cm, minimum height=1.2cm, inner sep=8pt] (Admin) {Administrador};
        \node[agente, right of=Admin, node distance=14cm, minimum width=4cm, minimum height=1.2cm, inner sep=8pt] (Sistema) {Sistema};

        % ==== 2. Almacén de Datos (BD) ====
        \node[almacen, below of=Admin, node distance=12cm, xshift=7cm, minimum width=6cm, inner sep=10pt] (BD) {BD EKIS (Publicidad)};

        % ==== 3. Procesos (RFs) ====
        \node[proceso, below of=Admin, node distance=5cm, xshift=-3cm, minimum width=4.5cm, inner sep=8pt] (RF31) {RF3.1: Añadir Anuncio};
        \node[proceso, below of=RF31, node distance=4cm, minimum width=4.5cm, inner sep=8pt] (RF32) {RF3.2: Añadir Característica};
        \node[proceso, right of=RF31, node distance=8cm, minimum width=4.5cm, inner sep=8pt] (RF33) {RF3.3: Eliminar Anuncio};
        \node[proceso, below of=RF33, node distance=4cm, minimum width=4.5cm, inner sep=8pt] (RF34) {RF3.4: Modificar Anuncio};
        \node[proceso, right of=RF34, node distance=8cm, minimum width=4.5cm, inner sep=8pt] (RF35) {RF3.5: Consultar Anuncio};

        % ==== 4. Flujos ====

        % % RF3.5 <-> Sistema
        \draw[flujo] (Sistema.south) -- node[midway, right, yshift=8pt] {RDE3.5} (RF35.north);
        \draw[flujo] (RF35.east) -- ++(1,0) |- node[pos=0.25, right, xshift=8pt] {RDS3.5} (Sistema.south east);

        % % RF3.1: Añadir Anuncio
        \draw[flujo] (Admin.south) -| ++(0,-2) -| node[pos=0.3, right, yshift=6pt] {RDE3.1} (RF31.north);
        % \draw[flujo, autoavoid] (RF31.east) to[out=0, in=180, looseness=1.2] node[midway, left, yshift=6pt] {RDW3.1} (BD.west);
        \draw[flujo, autoavoid] (RF31.east) -- ++(2,0) |- node[midway, left, yshift=6pt] {RDW3.1} (BD.west);

        % % RF3.2: Añadir Característica
        \draw[flujo, autoavoid] (Admin.south west) |- ++(-1,-1) -| node[pos=0.2, above, xshift=-8pt] {RDE3.2} (RF32.north west);
        \draw[flujo] (RF32.south) |- ++(0,-5) -| node[pos=0.4, above, xshift=-8pt] {RDW3.2} (BD.south);

        % % RF3.3: Eliminar Anuncio
        \draw[flujo] (Admin.south) |- ++(0,-2) -| node[pos=0.3, above, yshift=6pt] {RDE3.3} (RF33.north);
        % \draw[flujo] (RF33.east) to[out=0, in=0] node[pos=0.3, left, xshift=-6pt] {RDW3.3 / RDR3.3} (BD.north);
        \draw[flujo] (RF33.east)  -- ++(1,0) -- ++(0,-4)  -- node[pos=0.3, left, xshift=15pt, yshift=60pt] {RDW3.3 / RDR3.3} (BD.north);

        % % RF3.4: Modificar Anuncio
        %\draw[flujo] (Admin.south) |- ++(0,-2) -| node[pos=0.3, right, xshift=-1pt, yshift=-30pt] {RDE3.4} (RF34.west);
        \draw[flujo] (Admin.south) |- ++(0,-4) -| node[pos=0.6, above, xshift=-30pt, yshift=20pt] {RDE3.4} (RF34.north west);


        \draw[flujo] (RF34.south) |- ++(0,-1) -| node[pos=0.3, above, xshift=8pt] {RDW3.4} (BD.north);

        % % RF3.5: Consultar Anuncio
        \draw[flujo] (BD.south east) |- ++(1,-1) -| node[pos=0.7, right, xshift=8pt] {RDR3.5} (RF35.south);

    \end{tikzpicture}
    } % Fin del resizebox
    \caption{DFD1 - Subsistema de Publicidad }
    \label{fig:dfd1_pub}
\end{figure}


\begin{figure}[h!]
    \centering
    \caption{DFD Esquema Armazón (DFD0) del Sistema EKIS}
    \resizebox{\textwidth}{!}{%
    \begin{tikzpicture}[node distance=4cm, auto, thick]

        % ==== 1. Agentes Externos ====
        \node[agente, minimum width=3cm] (Admin) {Administrador};
        \node[agente, right of=Admin, node distance=5cm, yshift=2cm] (Usuario) {Usuario};
        \node[agente, right of=Usuario, node distance=5cm, yshift=0cm] (Visitante) {Visitante};
        \node[agente, below of=Usuario, node distance=9cm] (Sistema) {Sistema}; % Agente interno/disparador

        % ==== 2. Almacén de Datos (BD) ====
        \node[almacen, below of=Admin, node distance=12cm, xshift=7cm, minimum width=7cm, inner sep=10pt] (BD) {BD EKIS (Datos Centrales)};

        % % ==== 3. Procesos (Subsistemas - P0) ====
        % % Colocación de los subsistemas en el centro
        % \node[proceso, below of=Admin, node distance=3cm] (P3) {P3: Publicidad}; % RF3.X
        % \node[proceso, below of=P3, node distance=2.5cm, xshift=-1cm] (P2) {P2: Tendencias}; % RF2.X
        % \node[proceso, below of=P2, node distance=2.5cm, xshift=-1cm] (P1) {P1: Publicaciones}; % RF1.X
        % \node[proceso, right of=P1, node distance=6cm] (P4) {P4: Usuarios}; % RF4.X / RF5.X
        % \node[proceso, right of=P4, node distance=4.5cm] (P5) {P5: Mensajería Privada}; % RF5.X / RF5.X

        % % Reorganizar la posición de la BD para centrar los flujos
        % \node[almacen, below of=P1, node distance=4.5cm, xshift=5cm, minimum width=8cm, inner sep=10pt] (BD_EKIS) {BD EKIS (Datos Centrales)};

        % % Reposicionamiento final para un layout claro
        % \node[proceso, above of=P3, node distance=1.5cm, yshift=0cm] (P4_User) {P4: Usuarios};
        % \node[proceso, below of=P3, node distance=2.5cm] (P1_Pub) {P1: Publicaciones};
        % \node[proceso, right of=P1_Pub, node distance=4cm] (P2_Trend) {P2: Tendencias};
        % \node[proceso, below of=P1_Pub, node distance=2.5cm] (P5_Msg) {P5: Mensajería Privada};

        % % Ajuste final de posiciones de los agentes
        % \node[agente] (Admin_new) at (P3.north west) [anchor=south east, xshift=-0.5cm, yshift=2.5cm] {Administrador};
        % \node[agente] (Usuario_new) at (P2_Trend.north east) [anchor=south west, xshift=0.5cm, yshift=2.5cm] {Usuario};
        % \node[agente, below of=Usuario_new, node distance=2.5cm, xshift=1.5cm] (Visitante_new) {Visitante};
        % \node[agente, below of=P2_Trend, node distance=4.5cm, xshift=2cm] (Sistema_new) {Sistema};
        
        % % Colocación central de procesos
        % \node[proceso] (P4) at (0, 0) {P4: Usuarios};
        % \node[proceso] (P1) at (-4, -2) {P1: Publicaciones};
        % \node[proceso] (P3) at (4, -2) {P3: Publicidad};
        % \node[proceso] (P5) at (-4, -5) {P5: Mensajería Privada};
        % \node[proceso] (P2) at (4, -5) {P2: Tendencias};
        
        % \node[almacen] (BD_EKIS) at (0, -8) {BD EKIS (Datos Centrales)};

        % % Agentes externos
        % \node[agente] (Usuario) at (0, 3) {Usuario};
        % \node[agente] (Admin) at (4, 3) {Administrador};
        % \node[agente] (Visitante) at (-4, 3) {Visitante};
        % \node[agente] (Sistema) at (7, 0) {Sistema};

        % % ==== 4. Flujos: Agentes Externos <-> Procesos ====

        % % P4: Usuarios (Usuario, Visitante)
        % \draw[flujo] (Visitante.south) |- (P4.north west) node[midway, left] {RDE5.1 (Registro)}; [3]
        % \draw[flujo] (P4.north) -- (Usuario.south) node[midway, left] {RDS (Conf. Mod/Elim/Bloq/Amistad)}; [4-7]
        % \draw[flujo] (Usuario.south) -- (P4.north) node[midway, right] {RDE5.X (Gestión Usuario)}; [4-7]

        % % P1: Publicaciones (Usuario)
        % \draw[flujo] (Usuario.south west) -- (P1.north) node[midway, right] {RDE1.X (Crear/Mod/Like/Elim)}; [8-11]
        % \draw[flujo] (P1.north) -- (Usuario.south west) node[midway, left] {RDS1.X (Listado/Conf.)}; [8-10, 12, 13]

        % % P5: Mensajería Privada (Usuario)
        % \draw[flujo] (P5.north) -- (Usuario.south west) node[midway, left, xshift=-25pt] {RDS5.1 (Mensajes/Listado)}; [14, 15]
        % \draw[flujo] (Usuario.south west) |- (P5.north) node[midway, right, xshift=-25pt] {RDE5.X (Enviar/Eliminar)}; [14, 16]

        % % P3: Publicidad (Admin, Sistema)
        % \draw[flujo] (Admin.south) -- (P3.north) node[midway, left] {RDE3.X (Gestión Anuncios)}; [17-20]
        % \draw[flujo] (P3.north) -- (Admin.south) node[midway, right] {RDS3.X (Conf. Gestión)}; [17-20]
        % \draw[flujo] (Sistema.west) |- (P3.east) node[midway, above] {RDS3.5 (Datos Anuncio)}; [21]
        % \draw[flujo] (P3.east) -- (Sistema.west) node[midway, below] {RDE3.5 (Consulta)}; [21]

        % % P2: Tendencias (Admin, Usuario, Sistema)
        % \draw[flujo] (Usuario.south east) -- (P2.north) node[midway, left] {RDE2.4 (Consulta Listado)}; [22]
        % \draw[flujo] (P2.north) -- (Usuario.south east) node[midway, right] {RDS2.X (Tendencias Listadas)}; [22, 23]
        % \draw[flujo] (Admin.south) |- (P2.north east) node[midway, above] {RDE2.X (Asignar/Eliminar T.)}; [24, 25]
        % \draw[flujo] (Sistema.west) -- (P2.east) node[midway, above] {RDE2.1 (Notificación Hashtag)}; [26]

        % % ==== 5. Flujos: Procesos <-> BD (R/W) ====

        % P4 <-> BD
        % \draw[flujo] (P4.south) -- (BD_EKIS.north west) node[midway, left] {Datos Usuario (W5.X)}; [6, 7, 27-29]
        % \draw[flujo] (BD_EKIS.north west) -- (P4.south) node[midway, left, yshift=-10pt] {Datos Usuario (R)};

        % % P1 <-> BD
        % \draw[flujo] (P1.south) -- (BD_EKIS.north west) node[midway, left] {Datos Publicación (W1.X)}; [8-10, 13]
        % \draw[flujo] (BD_EKIS.north west) -- (P1.south) node[midway, right, xshift=-15pt] {Datos Publicación (R1.X)}; [9, 10, 12]

        % % P5 <-> BD
        % \draw[flujo] (P5.south) -- (BD_EKIS.west) node[midway, left] {Datos Mensajería (W5.X)}; [16, 30]
        % \draw[flujo] (BD_EKIS.west) -- (P5.south) node[midway, right, yshift=-10pt] {Datos Mensajería (R5.1)}; [15]

        % % P3 <-> BD
        % \draw[flujo] (P3.south) -- (BD_EKIS.north east) node[midway, right] {Datos Publicidad (W3.X)}; [17-19, 31]
        % \draw[flujo] (BD_EKIS.north east) -- (P3.south) node[midway, left, xshift=15pt] {Datos Publicidad (R3.5)}; [21]

        % % P2 <-> BD
        % \draw[flujo] (P2.south) -- (BD_EKIS.east) node[midway, right] {Datos Tendencias (W2.X)}; [24-26]
        % \draw[flujo] (BD_EKIS.east) -- (P2.south) node[midway, left, yshift=-10pt] {Datos Tendencias (R2.X)}; [22-24, 26, 32]

    \end{tikzpicture}
    }
\end{figure}


 \chapter{Usuarios}
 \begin{figure}[h!]
    \centering
    \resizebox{\textwidth}{!}{%
    \begin{tikzpicture}[node distance=4cm, auto]

        % ==== 1. Agentes Externos (fila superior) ====
        \node[agente, minimum width=4cm, minimum height=1.2cm, inner sep=8pt] (Visitante) at (-10,6) {Visitante};
        \node[agente, minimum width=4cm, minimum height=1.2cm, inner sep=8pt] (Usuario)   at ( 10,6) {Usuario};

        % ==== 2. Procesos (fila intermedia, muy separados) ====
        \node[proceso, minimum width=4.5cm, inner sep=8pt] (RF41) at (-10,2) {RF4.1: Crear};
        \node[proceso, minimum width=4.5cm, inner sep=8pt] (RF42) at ( -5,2) {RF4.2: Modificar};
        \node[proceso, minimum width=4.5cm, inner sep=8pt] (RF43) at (  0,2) {RF4.3: Eliminar};
        \node[proceso, minimum width=4.5cm, inner sep=8pt] (RF44) at (  5,2) {RF4.4: Bloquear};
        \node[proceso, minimum width=4.5cm, inner sep=8pt] (RF45) at ( 10,2) {RF4.5: A\~{n}adir Amigo};

        % ==== 3. Almacén de Datos (fila inferior) ====
        \node[almacen, minimum width=6cm, inner sep=10pt] (BD) at (0,-3.5) {EKIS (Usuarios y Relaciones)};

        % ==== 4. Flujos de datos ====

        % --- RF4.1: Crear Usuario (Visitante -> RF4.1 -> BD) ---
        \draw[flujo] (Visitante.south) -- node[midway,left] {RDE4.1} (RF41.north);
        \draw[flujo] (RF41.south) -- ++(0,-2) -| node[pos=0.0,below] {RDW4.1} (BD.west);

        % --- RF4.2: Modificar Usuario (Usuario -> RF4.2 -> BD) ---
        \draw[flujo] (Usuario.south) to[out=180,in=60] node[midway,above] {RDE4.2} (RF42.north);
        \draw[flujo] (RF42.south) -- ++(0,-1.2) -| node[pos=0.0,below] {RDW4.2} ($(BD.north west)+(1,0)$);

        % --- RF4.3: Eliminar Usuario (Usuario -> RF4.3 -> BD) ---
        \draw[flujo] (Usuario.south) to[out=200,in=60] node[midway,above] {RDE4.3} (RF43.north);
        \draw[flujo] (RF43.south) -- node[midway,right] {RDW4.3} (BD.north);

        % --- RF4.4: Bloquear Usuario (Usuario -> RF4.4 -> BD) ---
        \draw[flujo] (Usuario.south) to[out=220,in=60] node[midway,above] {RDE4.4} (RF44.north);
        \draw[flujo] (RF44.south) -- ++(0,-1.2) -| node[pos=0.0,below] {RDW4.4} ($(BD.north east)-(1,0)$);

        % --- RF4.5: A\~{n}adir Amigo (Usuario -> RF4.5 -> BD) ---
        \draw[flujo] (Usuario.south) -- node[midway,right] {RDE4.5} (RF45.north);
        \draw[flujo] (RF45.south) -- ++(0,-2) -| node[pos=0.0,below] {RDW4.5} (BD.east);

    \end{tikzpicture}
    } % Fin del resizebox
    \caption{DFD1 - Subsistema de Usuarios}
    \label{fig:dfd1_usuarios}
\end{figure}

 \chapter{Mensajería Privada}
 \begin{figure}[h!]
    \centering
    \shorthandoff{<>}
    \resizebox{\textwidth}{!}{
    \begin{tikzpicture}[node distance=2.2cm, auto]
    
        % ======================================================
        % === COLUMNA IZQUIERDA: AGENTE ========================
        % ======================================================
        \node[agente, minimum width=1cm, minimum height=1.2cm] (Usuario) {Usuario};

        % ======================================================
        % === COLUMNA CENTRAL: PROCESOS ========================
        % ======================================================
        
        \node[proceso, right of=Usuario, node distance = 8cm, minimum width=5cm] (P1) {RF5.1: Enviar mensaje};
        \node[proceso, below of=P1,      node distance = 3cm, minimum width=5cm] (P2) {RF5.2: Eliminar mensaje};
        \node[proceso, below of=P2,      node distance = 3cm, minimum width=5cm] (P3) {RF5.3: Listar usuarios};
        \node[proceso, below of=P3,      node distance = 3cm, minimum width=5cm] (P4) {RF5.4: Visualizar mensajes};
        \node[proceso, below of=P4,      node distance = 3cm, minimum width=5cm] (P5) {RF5.5: Des/Archivar usuario};
    
        % ======================================================
        % === COLUMNA DERECHA: ALMACENES ==================
        % ======================================================
        \node[almacen, right of=P3, node distance=10cm, minimum width=5cm] (Convers) {Conversaciones};
    
        % ======================================================
        % === FLUJOS DE DATOS ==================================
        % ======================================================
    
        % -- Usuario → Procesos --
        \draw[flujo] (Usuario) -- node[midway, above] {RDE5.1} (P1);
        \draw[flujo] ($(Usuario.south)+(+0.75, 0)$) |- node[pos=0.8, above] {RDE5.2} (P2);
        \draw[flujo] ($(Usuario.south)+(+0.45,0)$) |- node[pos=0.8, above] {RDE5.3} ($(P3.west)+(0, 0.1)$);
        \draw[flujo] ($(Usuario.south)+(-0.15,0)$) |- node[pos=0.8, above] {RDE5.4} ($(P4.west)+(0, 0.1)$);
        \draw[flujo] ($(Usuario.south)+(-0.75,0)$) |- node[pos=0.8, below] {RDE5.5} (P5);
    
        % -- Procesos → Usuario --
        \draw[flujo] ($(P3.west)+(0,-0.1)$) -| node[pos=0.2, below] {RDS5.1} ($(Usuario.south)+(+0.15,0)$);
        \draw[flujo] ($(P4.west)+(0,-0.1)$) -| node[pos=0.2, below] {RDS5.2} ($(Usuario.south)+(-0.45,0)$);

        % -- Procesos -> Almacenes --
        \draw[flujo] ($(P1.east)+(0,0.25)$) -| node[pos=0.2, above] {RDW5.1} ($(Convers.north)+(1.5,0)$);
        \draw[flujo] ($(P2.east)+(0,0.25)$) -| node[pos=0.2, above] {RDW5.2} ($(Convers.north)+(-0.5,0)$);
        \draw[flujo] ($(P4.east)+(0,-0.25)$) -| node[pos=0.2, below] {RDW5.3} ($(Convers.south)+(-0.5,0)$);
        \draw[flujo] ($(P5.east)+(0,-0.25)$) -| node[pos=0.2, below] {RDW5.4} ($(Convers.south)+(1.5,0)$);

        % -- Procesos <- Almacenes --
        \draw[flujo] ($(Convers.north)+(0.5,0)$) |- node[pos=0.8, below] {RDR5.1} ($(P1.east)+(0,-0.25)$);
        \draw[flujo] ($(Convers.north)+(-1.5,0)$) |- node[pos=0.8, below]{RDR5.2} ($(P2.east)+(0,-0.25)$);
        \draw[flujo] (Convers) -- node[pos=0.5, below] {RDR5.3} (P3);
        \draw[flujo] ($(Convers.south)+(-1.5,0)$) |- node[pos=0.8, above] {RDR5.4} ($(P4.east)+(0,0.25)$);
        \draw[flujo] ($(Convers.south)+(0.5,0)$) |- node[pos=0.8, above] {RDR5.5} ($(P5.east)+(0,0.25)$);

    \end{tikzpicture}
    }
    \caption{DFD1 — Subsistema de mensajería privada}
    \label{fig:dfd_mensajeria}
    \end{figure}
    

 \chapter{DFD0}
 A continuación, abordaremos el DFD0 de nuestra aplicación en la figura \ref{fig:dfd0} :

\begin{figure}[h!]
    \centering
    \shorthandoff{<>}
    \resizebox{\textwidth}{!}{
    \begin{tikzpicture}[node distance=2.2cm, auto]
    
        % ======================================================
        % ===================  AGENTES  ========================
        % ======================================================
        \node[agente, minimum width=1cm, minimum height=1.2cm] (Visitante) {Visitante};
        \node[agente, below of= Visitante    , node distance = 4cm, minimum width=1cm, minimum height=1.2cm] (Usuario) {Usuario};
        \node[agente, below of= Usuario      , node distance = 4cm, minimum width=1cm, minimum height=1.2cm] (Administrador) {Administrador};
        \node[agente, below of= Administrador, node distance = 4cm, minimum width=1cm, minimum height=1.2cm] (Sistema) {Sistema};
        
        % ======================================================
        % =================== PROCESOS ========================
        % ======================================================
        
        \node[proceso, right of=Visitante,     node distance = 8cm, minimum width=5cm] (S1) {Usuarios};
        \node[proceso, below of=S1,            node distance = 3cm, minimum width=5cm] (S2) {Publicaciones};
        \node[proceso, below of=S2,            node distance = 3cm, minimum width=5cm] (S3) {Tendencias};
        \node[proceso, below of=S3,            node distance = 3cm, minimum width=5cm] (S4) {Publicidad};
        \node[proceso, below of=S4,            node distance = 3cm, minimum width=5cm] (S5) {Mensajería privada};
    
        % ======================================================
        % ================== ALMACENES =========================
        % ======================================================
        \node[almacen, right of=S3, node distance=10cm, minimum height=4cm, minimum width=7cm] (EKIS) {EKIS};
    
        % ======================================================
        % === FLUJOS DE DATOS ==================================
        % ======================================================
    
        % -- Usuario → Procesos --
            
        % -- Procesos → Usuario --
        

        % -- Procesos -> Almacenes--
        
        \draw[->] (S1.east) -| ($(EKIS.north)+(1.5cm,0)$) node[pos=0.5, above] {Escritura usuarios};
        \draw[->] (S2.north east) -| ($(EKIS.north)+(-1cm,0)$) node[pos=0.5, above] {Escritura publicaciones};
        \draw[->] (S3.north east) -- ($(EKIS.west)+(0,0.5cm)$) node[pos=0.5, above] {Escritura Tendencias};
        \draw[->] (S4.north east) -| ($(EKIS.south)+(-2cm,0)$) node[pos=0.5, below] {Escritura Publicidad};
        \draw[->] (S5.north east) -| ($(EKIS.south)+(1cm,0)$) node[pos=0.5, below] {Escritura Mensajes};
        
        % -- Procesos <- Almacenes --
        \draw[<-] (S2.south east) -| ($(EKIS.north)+(-2cm,0)$) node[pos=0.5, above] {Escritura publicaciones};
        \draw[<-] (S3.south east) -- ($(EKIS.west)+(0,-0.5cm)$) node[pos=0.5, above] {Escritura Tendencias};
        \draw[<-] (S4.south east) -| ($(EKIS.south)+(-1cm,0)$) node[pos=0.5, below] {Escritura Publicidad};
        \draw[<-] (S5.south east) -| ($(EKIS.south)+(2cm,0)$) node[pos=0.5, below] {Escritura Mensajes};

    \end{tikzpicture}
    }
    \caption{DFD0 — Sistema EKIS}
    \label{fig:dfd0}
    \end{figure}
    

 \chapter{Caja Negra}
 \begin{figure}[h!]
    \centering
    \shorthandoff{<>}
    \resizebox{\textwidth}{!}{
    \begin{tikzpicture}[node distance=2.2cm, auto]
    
        % ======================================================
        % ===================  AGENTES  ========================
        % ======================================================
        \node[agente, minimum width=1cm, minimum height=1.2cm] (Visitante) {Visitante};
        \node[agente, below of= Visitante    , node distance = 3cm, minimum width=1cm, minimum height=2.4cm] (Usuario) {Usuario};
        \node[agente, below of= Usuario      , node distance = 3cm, minimum width=1cm, minimum height=1.2cm] (Administrador) {Administrador};
        \node[agente, below of= Administrador, node distance = 3cm, minimum width=1cm, minimum height=1.2cm] (Sistema) {Sistema};
        
        % ======================================================
        % =================== PROCESOS ========================
        % ======================================================
        
        \node[proceso, minimum width=5cm, minimum height=10.2cm] (EKIS) at ($(Visitante) + (10cm,-4.5cm)$) {Sistema EKIS};
    
        % ======================================================
        % === FLUJOS DE DATOS ==================================
        % ======================================================
    
        % -- Usuario → Procesos --
        \draw[->] (Visitante.east) -- node[pos=0.5, above, sloped]{Crear Usuario} ($(EKIS.west)+(0,4.71cm)$);

        \draw[->] ($(Usuario.east)+(0,1.2cm)$) -- node[pos=0.65, above, sloped]{Gestionar Usuarios} ($(EKIS.west)+(0,3.925cm)$);
        \draw[->] ($(Usuario.east)+(0,0.8cm)$) -- node[pos=0.65, above, sloped]{Gestionar Publicaciones} ($(EKIS.west)+(0,3.14cm)$);
        \draw[->] ($(Usuario.east)+(0,0.0cm)$) -- node[pos=0.65, above, sloped]{Gestionar Mensajes} ($(EKIS.west)+(0,1.57cm)$);
        \draw[->] ($(Usuario.east)+(0,-0.8cm)$) -- node[pos=0.65, above, sloped]{Peticion hashtag} ($(EKIS.west)+(0,-0.785cm)$);

        \draw[->] ($(Administrador.east)+(0,0.4cm)$) -- node[pos=0.5, above, sloped]{Gestionar Tendencias} ($(EKIS.west)+(0,-2.355cm)$);
        \draw[->] ($(Administrador.east)+(0,-0.4cm)$) -- node[pos=0.5, above, sloped]{Gestionar Publicidad} ($(EKIS.west)+(0,-3.14cm)$);

        \draw[->] ($(Sistema.east)+(0,0.4cm)$) -- node[pos=0.5, above, sloped]{Consultar Anuncio} ($(EKIS.west)+(0,-3.925cm)$);

        % -- Procesos → Usuario --
        \draw[<-] ($(Usuario.east)+(0,0.4cm)$) -- node[pos=0.65, above, sloped]{Salida Publicaciones} ($(EKIS.west)+(0,2.355cm)$);
        \draw[<-] ($(Usuario.east)+(0,-0.4cm)$) -- node[pos=0.65, above, sloped]{Salida Mensajes} ($(EKIS.west)+(0,0.785cm)$);
        \draw[<-] ($(Usuario.east)+(0,-1.2cm)$) -- node[pos=0.65, above, sloped]{Salida Tendencias} ($(EKIS.west)+(0,-1.57cm)$);

        \draw[<-] ($(Sistema.east)+(0,-0.4cm)$) -- node[pos=0.5, above, sloped]{Salida Anuncio} ($(EKIS.west)+(0,-4.71cm)$);


    \end{tikzpicture}
    }
    \caption{DFD0 — Sistema EKIS}
    \label{fig:caja_negra}
    \end{figure}
    En este apartado, especificamos el esquema de caja negra de nuestra aplicación, que se puede observar en la figura \ref{fig:caja_negra}. 
    Al igual que en la figura \ref{fig:dfd0}, los flujos de datos se han renombrado de igual manera para la simplicidad y mejor legebilidad del código.


 \chapter{Esquema externo}
 
\begin{figure}[h!]
    \centering
    \shorthandoff{<>}
    \resizebox{\textwidth}{!}{
    \begin{tikzpicture}[node distance=2.2cm, auto]
    
        % ======================================================
        % ===================  AGENTES  ========================
        % ======================================================
        \node[hexnode] (Visitante) {};
        \node[smallbox] at (Visitante.center) {Visitante};
        % ======================================================
        % =================== PROCESOS ========================
        % ======================================================
        
        \node[proceso, minimum width=5cm, minimum height=1.2cm] (EKIS) at ($(Visitante) + (10cm,0)$) {Sistema EKIS};
    
        % ======================================================
        % === FLUJOS DE DATOS ==================================
        % ======================================================
    
        % -- Usuario → Procesos --
        \draw[thick] (Visitante) -- (EKIS)

    \end{tikzpicture}
    }
    \caption{DFD0 — Sistema EKIS}
    \label{fig:mensajeria_esquema_externo}
    \end{figure}



% ----------------------- %
% BIBLIOGRAFÍA
% ----------------------- %

% Estilo de cita.
\bibliographystyle{apa-good}
%\renewcommand\bibname{Reference}

%[citestyle=numeric]

%\renewcommand\bibname{Bibliografía}
%

% Añadimos la bibliografía al índice
\phantomsection
\end{document}
